\chapter{現代の LLM アーキテクチャ}\label{ch:modern-arch}
\bookmap{3}

\begin{chapterlead}
\textbf{ここまで}: 注意機構を部品として、GPT-2 と同じ設計のミニGPT を組み立て、学習させ、文章を生成させました（\chapref{ch:transformer}）。\\[1mm]
\textbf{この章の問い}: 2019 年の GPT-2 から現在の LLM までに、設計のどこが、なぜ変わったのでしょうか。
KV キャッシュや MoE とは何でしょうか。\\[1mm]
\textbf{この章を終えると}:
\begin{itemize}
\item RMSNorm、RoPE、SwiGLU、GQA が「何に困って生まれたか」を説明し、PyTorch で実装できます。
\item KV キャッシュを自作し、生成が何倍速くなるかを測り、必要なメモリを式で見積もれます。
\item 混合エキスパート（MoE）の層を実装し、総パラメータとアクティブパラメータの違いを説明できます。
\item Llama 3、Qwen3、DeepSeek-V3、gpt-oss などの公開された設定を読み、パラメータ数とメモリを計算できます。
\end{itemize}
\end{chapterlead}

\begin{readingguide}
\item[\rgpre] 要るのは、\chapref{ch:attention}の注意の式と因果マスク、\chapref{ch:transformer}のブロックと \code{generate} です。初めて本書を読む人は、この章を後に回して\chapref{ch:pretraining}へ進めます。その場合も、KV キャッシュ（\secref{sec:ma-kvcache}の\figref{fig:ma-kvcache}まで）と MoE（\secref{sec:ma-moe}の最初の小節）だけは、先に読んでおくと安心です。
\item[\rgmust] \secref{sec:ma-diff}と、\secref{sec:ma-rope}の前半（2 次元の回転まで）：差分の一覧と、この章でいちばん大事な数学である「内積に位置が入るのは、位置の差を通してだけ」ということの導出。
\item[\rgmust] \secref{sec:ma-kvcache}、\secref{sec:ma-gqa}の前半、\secref{sec:ma-moe}（実装と実験を除く）：KV キャッシュ、メモリの式\eqref{eq:ma-kv-mem}、MoE。\chapref{ch:infra}と\chapref{ch:eval-serving}が、そのまま使います。
\item[\rgskim] \secref{sec:ma-rmsnorm}、\secref{sec:ma-swiglu}、\secref{sec:ma-bias}、\secref{sec:ma-assemble}、\secref{sec:ma-models}と、実装の小節：見出しの下に、押さえる結論を書きました。
\item[\rglater] FlashAttention（\secref{sec:ma-gqa}の最後）と\secref{sec:ma-beyond}：\chapref{ch:gpu}で高速化を学ぶときや、読み物として。
\end{readingguide}

前の章の最後で、本書のミニGPT は 2019 年の GPT-2 と同じ設計だと述べました。
では、2025 年や 2026 年に公開された LLM の中身は、どれほど違うのでしょうか。
意外なことに、骨組みはほとんど同じです。
埋め込み、ブロックの積み重ね、注意と FFN、残差接続、出力ヘッド。
この流れは今も変わっていません。
FFN（feed-forward network）は、第 10 章で MLP と呼んだ部品のことです。
現在の LLM の論文や設定ファイルではこの名前がよく使われるので、本章ではこちらで呼びます。

変わったのは、ブロックの中の部品です。
正規化の式、位置の教え方、FFN の形、注意のヘッドの持ち方。
一つひとつは小さな変更ですが、どれにも「大きくしたときに困ったこと」への答えという理由があります。
もう一つ大きく変わったのが、\textbf{推論の効率}への配慮です。
LLM は学習の後に、何億回も文章を生成するのに使われます。
生成を速く、安くするために、生成の計算のやり方が工夫され、それに合わせてモデルの設計も変わりました。
その効果は、この章で数字として確かめます。
過去の計算を取っておく「KV キャッシュ」を自作すると、本書の小さなモデルでも、長い生成が約 8 倍速くなります。
注意のヘッドの持ち方を変えるだけで、会話 1 本ごとに取っておく KV キャッシュのメモリが 4 分の 1 に減ることも計算します。
2026 年に公開されたモデルが、何層に何個の注意のヘッドを重ねているのかも、\secref{sec:ma-models}で本書のミニGPT と並べて見ます。

この章では、GPT-2 型から現代の標準形への差分を一つずつ取り上げます。
それぞれについて「何に困っていたか → どう変えたか → コードではどう書くか」を確かめます。
最後に、それらをすべて入れた「現代版ミニGPT」を作り、基本版と同じ条件で学習させて比べます。

\section{GPT-2 から何が変わったか：差分の一覧}\label{sec:ma-diff}

細部に入る前に、変更点の全体を一枚の表で眺めておきましょう。
この節の問いは、「GPT-2 と現代の LLM は、部品ごとに見ると何が違うのか」です。

\tabref{tab:ma-diff}が差分の一覧です。
左の列が本書のミニGPT と GPT-2 の部品、右の列が Llama 3 や Qwen3 など現在の公開モデルの多くが採用している部品です。
「現代の標準」は \cite{touvron2023llama} の LLaMA がまとめた形が出発点で、その後のモデルの多くがこれを受け継いでいます。

\begin{table}[tbp]
\centering
\caption{GPT-2 型（第 10 章のミニGPT）と現代の標準的な LLM の差分。上から順に本章で一つずつ扱う。罫線で、ねらいの同じ行をまとめた。「生成のやり方」の行だけは部品の差ではなく計算のやり方の差で、GPT-2 型でも KV キャッシュは使える。$4d$ などは FFN の中間の次元。「密」は、すべてのトークンが同じ 1 つの FFN を通ること（\ref{sec:ma-moe}節）。}
\label{tab:ma-diff}
\footnotesize
\setlength{\tabcolsep}{3pt}
\begin{tabular}{@{}lllll@{}}
\toprule
部品 & GPT-2 型 & 現代の標準 & ねらい & 本章 \\
\midrule
正規化 & LayerNorm（部品の前） & RMSNorm（部品の前） & 学習を安定に、安く & \ref{sec:ma-rmsnorm}節 \\
位置の情報 & 学習型の位置埋め込み & RoPE（回転） & 相対位置、長い文脈 & \ref{sec:ma-rope}節 \\
FFN & GELU の 2 層 MLP（$4d$） & SwiGLU（約 $\tfrac{8}{3}d$〜$3.5d$） & 学習をうまく & \ref{sec:ma-swiglu}節 \\
バイアス & あり & なし & 部品を減らす & \ref{sec:ma-bias}節 \\
\midrule
埋め込みと出力 & 重み共有 & 小さいと共有、大きいと別 & 大きさで選ぶ & \ref{sec:ma-bias}節 \\
\midrule
生成のやり方 & 毎回すべて計算し直す\textsuperscript{*} & KV キャッシュ & 生成を速く & \ref{sec:ma-kvcache}節 \\
注意のヘッド & マルチヘッド（MHA） & GQA（キー・バリューを共有） & キャッシュを小さく & \ref{sec:ma-gqa}節 \\
\midrule
FFN の規模 & 1 つ（密） & 大きいモデルは MoE & 計算を増やさず大きく & \ref{sec:ma-moe}節 \\
\bottomrule
\multicolumn{5}{@{}l}{\textsuperscript{*}第 10 章の \code{generate} の書き方。GPT-2 型のモデルでも KV キャッシュは使える。}
\end{tabular}
\end{table}

表は、ねらいごとに 4 つのまとまりに分けてあります。
上の 4 行は、\textbf{学習をうまく、安くする}ための変更です（RoPE には、長い文脈を扱いやすくするねらいもあります）。
5 行目の重み共有は、パラメータを節約するかどうかの選択で、モデルの大きさによって答えが分かれます。
次の 2 行は、\textbf{生成を速くする}ための工夫です。
ただし性質が違います。
KV キャッシュは生成の計算のやり方の工夫で、モデルの部品（重みや式）は何も変えません。
GQA は、そのキャッシュを小さくするための、部品の変更です。
最後の MoE は、\textbf{1 トークンあたりの計算量をほとんど増やさずに、パラメータを増やす}ための変更です。

\figref{fig:ma-block}は、1 つのブロックを左右に並べて比べたものです。
データの流れ（残差ストリームに注意と FFN が足し込む）は同じで、箱の中身だけが入れ替わっていることを確かめてください。

\begin{figure}[tbp]
\centering
\begin{tikzpicture}[x=1mm,y=1mm,
  b/.style={box,font=\footnotesize,minimum width=30mm,minimum height=6mm},
  n/.style={graybox,font=\footnotesize,minimum width=30mm,minimum height=5.5mm}]
\foreach \sx/\ttl/\fa/\fb/\ta/\tb/\tc/\td in {
  0/{GPT-2 型}/figblue/bookblue/{LayerNorm}/{注意（MHA）}/{LayerNorm}/{FFN（MLP, GELU, $4d$）},
  64/{現代の標準}/figgreen/bookgreen/{RMSNorm}/{注意（GQA＋RoPE）}/{RMSNorm}/{SwiGLU（または MoE）}} {
\begin{scope}[xshift=\sx mm]
\node[font=\small\bfseries] at (14,64) {\ttl};
\draw[arr,very thick,bookblue!70] (0,-4) -- (0,57);
\node[lbl,anchor=north,align=center] at (0,-4.5) {前のブロックから};
\node[lbl,anchor=south] at (0,57.5) {次のブロックへ};
\node[lbl,rotate=90,anchor=south] at (-1.5,36) {残差ストリーム};
\node[n,fill=\fa,draw=\fb] (a1) at (20,8) {\ta};
\node[b,fill=\fa,draw=\fb] (a2) at (20,17) {\tb};
\node[op] (p1) at (0,24) {$+$};
\node[n,fill=\fa,draw=\fb] (a3) at (20,33) {\tc};
\node[b,fill=\fa,draw=\fb,minimum width=34mm] (a4) at (20,42) {\td};
\node[op] (p2) at (0,49) {$+$};
\draw[arr] (0,2) -| (a1.south);
\draw[arr] (a1) -- (a2); \draw[arr] (a2.north) |- (p1);
\draw[arr] (0,27) -| (a3.south);
\draw[arr] (a3) -- (a4); \draw[arr] (a4.north) |- (p2);
\end{scope}}
\end{tikzpicture}
\caption{1 つのブロックの比較。左が GPT-2 型（第 10 章）、右が現代の標準形（緑の箱が入れ替わった部品）。太い縦の矢印が残差ストリームで、各位置のベクトルは箱を通らずにそのまま上へ進む。注意と FFN は、横に分かれた枝で値を読み出し、出力を「＋」で足し込む。この流れは左右で同じ。}
\label{fig:ma-block}
\end{figure}

この節で分かったのは、現代の LLM が GPT-2 と「同じ骨組みに、別の部品」を入れたものだということです。
次の節から、表の上から順に部品を交換していきます。

\section{RMSNorm：平均を引くのをやめる}\label{sec:ma-rmsnorm}
\skimread{押さえる結論は「RMSNorm は、LayerNorm から『平均を引く』と『ずらし $\bm\beta$』を取り除き、大きさだけをそろえる軽い正規化」です。}

最初の差分は正規化です。
レイヤー正規化（LayerNorm、\chapref{ch:transformer}）は、各位置のベクトルを平均 0・分散 1 にそろえ、深いネットワークの学習を安定させる部品でした。
この節の問いは、「LayerNorm の計算のうち、本当に必要なのはどこか」です。

\subsection*{LayerNorm の 2 つの仕事}

LayerNorm は、$d$ 次元のベクトル $\vx$ に対して次の計算をしました。
\begin{equation}
\mathrm{LayerNorm}(\vx)=\bm{\gamma}\odot\frac{\vx-\mu}{\sqrt{\sigma^2+\epsilon}}+\bm{\beta},\qquad
\mu=\frac1d\sum_{i=1}^d x_i,\quad \sigma^2=\frac1d\sum_{i=1}^d(x_i-\mu)^2 .
\label{eq:ma-layernorm}
\end{equation}
ここで $\odot$ は成分ごとの積、$\epsilon$ は 0 で割らないための小さな数です。
$\bm\gamma$ は成分ごとの倍率、$\bm\beta$ は成分ごとのずらしで、どちらも学習するパラメータです（\chapref{ch:transformer}）。
この式は 2 つの仕事をしています。
一つは平均 $\mu$ を引いて\textbf{中心をそろえる}こと、もう一つは標準偏差で割って\textbf{大きさをそろえる}ことです。

Zhang と Sennrich\cite{zhang2019root}は、学習を安定させているのは主に後者だと考えました。
値が層を通るたびに大きくなったり小さくなったりするのを止められれば、中心は多少ずれていても困らない、という見立てです。
それなら、平均を計算して引く手間を省けます。
そこで生まれたのが、\term{RMSNorm}{root mean square normalization}です。

\begin{definition}[RMSNorm]
$d$ 次元のベクトル $\vx$ の 2 乗平均平方根（root mean square）を $\mathrm{RMS}(\vx)=\sqrt{\frac1d\sum_{i=1}^d x_i^2+\epsilon}$ とします。RMSNorm は
\begin{equation}
\mathrm{RMSNorm}(\vx)=\bm{\gamma}\odot\frac{\vx}{\mathrm{RMS}(\vx)}
\label{eq:ma-rmsnorm}
\end{equation}
です。学習するパラメータは倍率 $\bm{\gamma}\in\R^d$ だけで、ずらし $\bm{\beta}$ はありません。
\end{definition}

\begin{example}[3 次元の手計算]
$\vx=(1,2,3)$、$\bm{\gamma}=(1,1,1)$、$\epsilon=0$ とします。
2 乗の平均は $(1+4+9)/3=14/3\approx4.667$ なので、$\mathrm{RMS}(\vx)=\sqrt{4.667}\approx2.160$ です。
したがって、各成分を 2.160 で割って $\mathrm{RMSNorm}(\vx)\approx(1/2.160,\,2/2.160,\,3/2.160)=(0.463,\,0.926,\,1.389)$ になります。
一方 LayerNorm では、まず平均 $\mu=(1+2+3)/3=2$ を引いて $(-1,0,1)$ にします。
分散は $\bigl((-1)^2+0^2+1^2\bigr)/3=2/3$ なので、標準偏差 $\sqrt{2/3}\approx0.816$ で割って $(-1.225,\,0,\,1.225)$ です。
RMSNorm の出力は平均が 0 ではありませんが、2 乗平均はちょうど 1 になっています（$0.463^2+0.926^2+1.389^2\approx3.0$、3 で割って 1）。
\end{example}

\begin{figure}[tbp]
\centering
\begin{tikzpicture}
\begin{axis}[width=0.7\linewidth,height=42mm,ybar,bar width=5.5mm,ymin=-1.9,ymax=3.8,
  xmin=0.5,xmax=3.5,xtick={1,2,3},xticklabels={第 1 成分,第 2 成分,第 3 成分},ylabel={値},
  grid=none,ymajorgrids,extra y ticks={0},extra y tick labels={},extra y tick style={grid style={black!55}},
  legend pos=outer north east,legend style={font=\scriptsize},legend cell align=left,
  nodes near coords,nodes near coords style={font=\tiny,text=black,/pgf/number format/.cd,fixed,precision=3}]
\addplot[fill=bookgray!60,draw=none] coordinates {(1,1) (2,2) (3,3)};
\addlegendentry{入力 $\vx$}
\addplot[fill=bookblue,draw=none] coordinates {(1,-1.225) (2,0) (3,1.225)};
\addlegendentry{LayerNorm}
\addplot[fill=bookaccent,draw=none] coordinates {(1,0.463) (2,0.926) (3,1.389)};
\addlegendentry{RMSNorm}
\end{axis}
\end{tikzpicture}
\caption{$\vx=(1,2,3)$ を 2 つの正規化に通した結果（上の例の数値）。LayerNorm（青）は平均を引くので、棒が 0 の上下に分かれる。RMSNorm（赤）は全体を同じ倍率 $1/2.160$ で縮めるだけなので、$1:2:3$ の比と符号がそのまま残る。}
\label{fig:ma-norm}
\end{figure}

\figref{fig:ma-norm}は、この例を棒グラフにしたものです。
青の LayerNorm は「中心を 0 に動かしてから縮める」、赤の RMSNorm は「その場で縮めるだけ」であることを見てください。

つまり RMSNorm は、倍率 $\bm\gamma$ を掛ける前の段階では「大きさだけをそろえ、向きはそのまま残す」正規化です
（$\bm\gamma$ の成分が互いに違えば、成分ごとに倍率が違うので、向きも少し変わります）。
2 乗平均が 1 ということは 2 乗の和が $d$、つまりベクトルの長さが $\sqrt d$ だということです（上の例では長さ $\sqrt3$）。
そこで、$\bm\gamma$ がすべて 1 なら、RMSNorm はベクトルを原点からの距離 $\sqrt{d}$ の球面に乗せる操作だと考えることもできます。
計算は平均を求める 1 回分だけ軽くなり、パラメータは半分になります。
LayerNorm は $\bm\gamma$ と $\bm\beta$ で $2d$ 個、RMSNorm は $\bm\gamma$ の $d$ 個だけだからです。
Zhang と Sennrich の実験では、精度を落とさずに正規化の計算時間を減らせました\cite{zhang2019root}。
ただし、正規化はモデル全体の計算のごく一部です。
パラメータも、本章の現代版ミニGPT では全体の約 125 万個のうち 1,152 個が減るだけです（\secref{sec:ma-assemble}）。
得は小さいのですが、精度をほとんど失わずに済むので、採用が広がりました。
LLaMA\cite{touvron2023llama}がこれを採用して以来、ほぼすべての公開 LLM が RMSNorm を使っています。

\subsection*{コード}

\snippet{RMSNorm}{code:ma-rmsnorm}{ch11/modern_gpt.py}{rmsnorm}

5 行目の \code{weight} が $\bm{\gamma}$ で、初期値はすべて 1 です。
8 行目が $\mathrm{RMS}(\vx)$ の計算で、\code{x.pow(2).mean(dim=-1)} が最後の軸（$d$ 個の成分）の 2 乗平均、\code{.add(self.eps).sqrt()} が $\epsilon$ を足して平方根をとる部分です。
\code{keepdim=True} で形を $(B,T,1)$ に保つので、9 行目の割り算はブロードキャスト（\chapref{ch:pytorch}）で各位置のベクトル全体に効きます。
PyTorch にも同じ働きの \code{nn.RMSNorm} がありますが、本章では式との対応を見るために自分で書きました。

正規化を置く場所は、GPT-2 と同じく注意と FFN の\textbf{前}です（\figref{fig:ma-block}）。
置き場所は変わっていません。
正規化の種類を問わない言い方として、Pre-Norm（事前正規化）という名前もよく使われます（Pre-LN は、LayerNorm を部品の前に置く場合の名前です）。
Gemma のように、部品の前と後の両方に RMSNorm を置くモデルもあります。

この節で分かったのは、RMSNorm が LayerNorm から「平均を引く」と「ずらし $\bm\beta$」を取り除いた、より軽い正規化だということです。

\section{RoPE：位置を「回転」で教える}\label{sec:ma-rope}

次は位置の情報です。
この節は本章でいちばん数学が多い節ですが、使うのは 2 次元の回転行列と三角関数の加法定理だけです。
この節の問いは、「学習型の位置埋め込みの何が困るのか、そして回転でどう解決するのか」です。

\subsection*{何に困っているか}

GPT-2 型のモデルは、位置 $t$ ごとに学習したベクトル $\vp_t$ をトークンの埋め込みに足していました（\chapref{ch:transformer}）。
この方法には 2 つの弱点があります。

一つ目は、\textbf{表にない位置を扱えない}ことです。
位置の表 $\mP$ の行数は文脈長の上限 $T_{\max}$ で決まります。
本書の基盤モデルなら 128 行です。
129 番目の位置のベクトルは学習されていないので、文脈長を後から伸ばすことができません。

二つ目は、\textbf{本当に知りたいのは相対的な位置}だということです。
「直前のトークン」「3 つ前のトークン」という関係は、文のどこに現れても同じ意味を持ちます。
ところが絶対位置の埋め込みでは、位置 5 と位置 2 の関係も、位置 105 と位置 102 の関係も、別々のベクトルの組み合わせとして学ぶしかありません。

そこで欲しいのは、「クエリとキーの内積に、2 つのトークンの位置の情報が\textbf{位置の差}の形でだけ入る」仕組みです。
内積は、これまでどおりクエリとキーの中身（どのトークンか）で決まります。
そのうえで、位置については「何番目か」ではなく「何個離れているか」だけが効く、という形です。
内積は注意のスコア（\chapref{ch:attention}）なので、そうなれば、注意のスコアは絶対位置ではなく相対位置を反映します。

\subsection*{2 次元の回転から始める}

まず、クエリとキーが 2 次元のベクトルだとしましょう。
平面のベクトルを角度 $\phi$ だけ回す行列は
\begin{equation}
\mR(\phi)=\begin{pmatrix}\cos\phi & -\sin\phi\\ \sin\phi & \cos\phi\end{pmatrix}
\label{eq:ma-rot}
\end{equation}
です。
これを使って、位置 $m$ のクエリ $\vq$ を角度 $m\omega$ だけ、位置 $n$ のキー $\vk$ を角度 $n\omega$ だけ回します。
$\omega$ は回る速さを決める定数です。
回した後の内積を計算してみましょう。
\begin{align}
\bigl(\mR(m\omega)\vq\bigr)^\T\bigl(\mR(n\omega)\vk\bigr)
&=\vq^\T\mR(m\omega)^\T\mR(n\omega)\vk && \text{転置の性質 $(\mA\vx)^\T=\vx^\T\mA^\T$}\nonumber\\
&=\vq^\T\mR(-m\omega)\mR(n\omega)\vk && \text{回転の転置は逆回転}\nonumber\\
&=\vq^\T\mR\bigl((n-m)\omega\bigr)\vk . && \text{回転を続けると角度が足される}
\label{eq:ma-rope-rel}
\end{align}
2 行目は、$\mR(\phi)^\T$ の成分を見ると $\cos(-\phi)=\cos\phi$、$\sin(-\phi)=-\sin\phi$ より $\mR(-\phi)$ に等しいことから分かります。
3 行目は、三角関数の加法定理から $\mR(\alpha)\mR(\beta)=\mR(\alpha+\beta)$ となることを使いました。
たとえば積の左上の成分は $\cos\alpha\cos\beta-\sin\alpha\sin\beta=\cos(\alpha+\beta)$ です（ほかの成分も同様）。
結果には $m$ と $n$ が\textbf{差 $n-m$ の形でしか現れません}。
これが求めていた性質です。
$\vq$ と $\vk$ はそのまま残っているので、スコアは中身にもよります。
以下の図とコードの出力では、差を $m-n$（クエリの位置 $-$ キーの位置）で表示します。
式の $n-m$ とは符号が逆なだけです。
しかも回転はベクトルの長さを変えないので、$\vq$ と $\vk$ の大きさの情報はそのまま残ります。

\begin{example}[角度が見える数値例]
$\omega=1$（ラジアン）、$\vq=\vk=(1,0)$ とします。
式\eqref{eq:ma-rope-rel}の最後の形 $\vq^\T\mR(\phi)\vk$ にこの $\vq,\vk$ を入れると、残るのは $\mR(\phi)$ の左上の成分 $\cos\phi$ だけです（$\phi=(n-m)\omega$）。
同じ位置 $m=n=3$ なら、両方とも同じ角度だけ回るので $\phi=0$ で、内積は $\cos 0=1$ のままです。
$m=5,\,n=3$ なら $\phi=(3-5)\times1=-2$ で、内積は $\cos(-2)=\cos 2\approx-0.416$ です。
$m=7,\,n=5$ でも差は同じなので、やはり $-0.416$ になります。
\end{example}

\begin{figure}[tbp]
\centering
\begin{tikzpicture}[scale=1.1]
\foreach \cx/\m/\n/\ttl in {0/5/3/{$m=5,\ n=3$}, 4.9/7/5/{$m=7,\ n=5$}} {
  \begin{scope}[xshift=\cx cm]
  \draw[black!25] (0,0) circle (1.3);
  \draw[->,black!40] (-1.6,0) -- (1.6,0) node[right,font=\scriptsize,text=black!60] {第 1 成分};
  \draw[->,black!40] (0,-1.6) -- (0,1.6) node[above,font=\scriptsize,text=black!60] {第 2 成分};
  \pgfmathsetmacro{\aq}{\m*57.2958}\pgfmathsetmacro{\ak}{\n*57.2958}
  \draw[->,very thick,black!45] (0,0) -- (1.3,0);
  \node[font=\tiny,text=black!60,anchor=north west,inner sep=1pt] at (1.3,-0.08) {回す前 $(1,0)$};
  \draw[->,bookblue!60] plot[domain=0:\aq,samples=80,smooth] ({\x}:{0.95+0.12*\x/360});
  \node[font=\tiny,text=bookblue,fill=white,inner sep=0.5pt] at ({\aq*0.40}:{1.12+0.05*\aq/360}) {$\vq$ を \m\,rad};
  \draw[->,bookaccent!60] plot[domain=0:\ak,samples=60,smooth] ({\x}:{0.72});
  \node[font=\tiny,text=bookaccent,fill=white,inner sep=0.5pt] at ({\ak*0.55}:0.55) {$\vk$ を \n\,rad};
  \draw[->,very thick,bookblue] (0,0) -- (\aq:1.3) node[pos=1.18,font=\footnotesize] {$\vq$};
  \draw[->,very thick,bookaccent] (0,0) -- (\ak:1.3) node[pos=1.18,font=\footnotesize] {$\vk$};
  \draw[bookgreen,very thick] (\ak:0.36) arc[start angle=\ak,end angle=\aq,radius=0.36];
  \node[font=\tiny,text=bookgreen!70!black,fill=white,inner sep=0.5pt] at ({(\ak+\aq)/2-19*(\m==7)}:{0.52+0.1*(\m==7)}) {2\,rad};
  \node[font=\footnotesize] at (0,-1.95) {\ttl};
  \end{scope}
}
\end{tikzpicture}
\caption{2 次元の RoPE（$\omega=1$）。回す前のクエリ $\vq$ とキー $\vk$ は、ともに $(1,0)$（灰色の矢印）。$\vq$ を $m$ ラジアン、$\vk$ を $n$ ラジアン、反時計回りに回す（細い弧）。右の $\vq$ は 7 ラジアンで、1 周（約 6.28 ラジアン）を約 0.72 ラジアン超えている。左右で向きは違うが、間の角度（緑の弧）はどちらも $m-n=2$ ラジアンなので、内積 $\cos2$ は同じになる。}
\label{fig:ma-rope2d}
\end{figure}

\figref{fig:ma-rope2d}は、この 2 つの場合を図にしたものです。
左と右で $\vq$ と $\vk$ の向きはまったく違いますが、2 本の間の角度（緑の弧）は同じです。
内積は 2 本の長さと間の角度だけで決まるので、同じ値になります。

このように、クエリとキーを位置に比例した角度だけ回してから内積をとる方法を、\term{RoPE}{rotary position embedding}（回転位置埋め込み）\cite{su2024roformer}と呼びます。
注意のスコアに位置が入るのは差 $m-n$ を通してだけになり、角度は式で計算できるので、学習する位置の表も要りません。
もちろん、スコアは $\vq$ と $\vk$ の中身にもよります。
変わったのは、位置の入り方です。
残りの 2 つの小節は、これを実際のヘッドの次元に広げ、コードにする話です。

\subsection*{$d$ 次元への拡張}
\skimread{押さえる結論は「実際の RoPE は、ヘッドの成分を 2 つずつの組に分け、組ごとに違う速さで回す。速い組は近くの位置を、遅い組は遠くの位置を見分ける」です。}

実際のヘッドの次元 $\dhead$ は 64 や 128 です。
実際の RoPE\cite{su2024roformer}は、$\dhead$ 個の成分を 2 つずつの組に分け、組ごとに別の速さで回します。
組の番号 $i$ は、コードと同じく 0 から数えます。
第 $i$ 組（$i=0,1,\dots,\dhead/2-1$）の速さは
\begin{equation}
\omega_i=\theta^{-2i/\dhead},\qquad \theta=10000
\label{eq:ma-rope-freq}
\end{equation}
です。
$\theta=10000$ と指数の形は、Transformer を最初に提案した論文\cite{vaswani2017attention}の正弦波の位置エンコーディング（\chapref{ch:transformer}）と同じ慣習の値です。
$i=0$ の組は $\omega_0=1$ で、1 トークン進むごとに 1 ラジアン回ります。
$i$ が大きくなるほどゆっくり回り、最後の組はほとんど回りません。
たとえば $\dhead=64$ のまん中の組（$i=16$）は $\omega_{16}=10000^{-1/2}=0.01$ で、1 周に $2\pi/0.01\approx628$ トークンかかります。
まとめると、RoPE は、各ヘッドのクエリとキーを、位置 $m$ に応じてこの組ごとの角度 $m\omega_i$ だけ回してから内積をとる方法です。
バリューは回しません。
各組で式\eqref{eq:ma-rope-rel}が成り立つので、組の内積の和である全体の内積も、位置については差だけで決まります。
\figref{fig:ma-rope-pairs}は、$\dhead=8$ の小さな例で、4 つの組が位置ごとにどれだけ回るかを並べたものです。
1 行目の速い組は 1 位置ごとに 1 ラジアン回り、4 行目の遅い組はほとんど動きません。

\begin{figure}[tbp]
\centering
\begin{tikzpicture}[x=1mm,y=1mm]
\foreach \i/\w/\lab in {0/1/{1},1/0.1/{0.1},2/0.01/{0.01},3/0.001/{0.001}} {
  \pgfmathtruncatemacro{\ca}{2*\i+1}\pgfmathtruncatemacro{\cb}{2*\i+2}
  \node[font=\scriptsize,anchor=east] at (-2,-\i*11) {組 \i：$(x_{\ca},x_{\cb})$、$\omega_{\i}=\lab$};
  \foreach \m in {0,...,4} {
    \pgfmathsetmacro{\ang}{\m*\w*57.2958}
    \begin{scope}[shift={(8+\m*14,-\i*11)}]
    \draw[black!25] (0,0) circle (4);
    \draw[black!20] (0,0) -- (4,0);
    \draw[->,thick,bookblue] (0,0) -- (\ang:4);
    \end{scope}
  }
}
\foreach \m in {0,...,4} {
  \node[font=\scriptsize] at (8+\m*14,7.5) {$m=\m$};
}
\node[font=\scriptsize,anchor=east] at (-2,7.5) {位置};
\end{tikzpicture}
\caption{$\dhead=8$ の RoPE で、クエリの 8 成分を 4 つの組に分け、組ごとに角度 $m\omega_i$ だけ回す様子（式\eqref{eq:ma-rope-freq}で $\theta=10000$）。1 つの円が 1 つの組で、矢印は回した後の向き、灰色の線は回す前の向き。組 0 は 1 位置ごとに 1 ラジアン（約 57°）回り、組 1 は 0.1 ラジアン、組 2 と組 3 は、この範囲ではほとんど回らない。}
\label{fig:ma-rope-pairs}
\end{figure}

なぜ速さを変えるのでしょうか。
時計の針にたとえると分かりやすくなります。
秒針（速い組）は近い位置の細かな違いを区別できますが、60 秒で一周してしまうので遠い位置は区別できません。
時針（遅い組）は細かな違いは苦手ですが、長い時間を区別できます。
いろいろな速さの針を並べることで、近い距離も遠い距離も表せるのです。
RoPE の数で言えば、$\omega_0=1$ の組は約 6.28 トークン（$2\pi$）で 1 周するので、差 0 と差 6 はほぼ同じ角度になり、この組だけでは見分けにくくなります。
一方、$\dhead=64$ の第 8 組（$\omega_8=10000^{-1/4}=0.1$、1 周に約 63 トークン）では、差 6 は 0.6 ラジアンなので、差 0 と見分けられます。
このたとえで対応させているのは、「速い組は近い位置の違いを、遅い組は遠い位置の違いを見分けやすい」という点だけです。
時計の針と違って、モデルがどの組をどう使うかは学習で決まります。

\begin{figure}[tbp]
\centering
\begin{tikzpicture}
\begin{axis}[width=0.48\linewidth,height=0.42\linewidth,
  xlabel={組の番号 $i$},ylabel={1 周の長さ（トークン）},ymode=log,
  title={\footnotesize (a) 組ごとの周期 $2\pi/\omega_i$},title style={yshift=-1mm}]
\addplot[bookblue,mark=*,mark size=1pt] table[col sep=comma,x=i,y=period]{code/ch11/out/rope_freq.csv};
\end{axis}
\end{tikzpicture}\hfill
\begin{tikzpicture}
\begin{axis}[width=0.48\linewidth,height=0.42\linewidth,
  xlabel={位置の差 $m-n$},ylabel={内積 / 長さの積},ymin=-0.45,ymax=1.08,
  legend pos=north east,legend style={font=\tiny,row sep=-2pt},legend cell align=left,
  title={\footnotesize (b) $\vq,\vk$ の選び方と内積},title style={yshift=-1mm}]
\addplot[bookaccent,line width=1pt] table[col sep=comma,x=dist,y=score]{code/ch11/out/rope_decay.csv};
\addlegendentry{同じ向き（全成分 1）}
\addplot[bookorange!90!black,densely dashed,thick,mark=o,mark size=1.5pt,mark repeat=12,mark options={solid,fill=white}] table[col sep=comma,x=dist,y=same]{code/ch11/out/sweep_rope_decay.csv};
\addlegendentry{同じ向き（ランダム）}
\addplot[bookblue,thick] table[col sep=comma,x=dist,y=indep_mean]{code/ch11/out/sweep_rope_decay.csv};
\addlegendentry{無関係（200 組平均）}
\addplot[bookgray,thin] table[col sep=comma,x=dist,y=indep_one]{code/ch11/out/sweep_rope_decay.csv};
\addlegendentry{無関係（1 組）}
\end{axis}
\end{tikzpicture}
\caption{RoPE の性質（$\dhead=64$、$\theta=10000$）。(a) 最初の組は約 6 トークンで 1 周し、最後の組は約 4.7 万トークンかかる。縦軸は対数。(b) $\vq$ と $\vk$ を位置の差だけ離して回したときの内積を、2 本の長さの積で割った値（2 本のなす角の cos）。$\vq$ と $\vk$ が同じ向きのとき（赤: すべての成分が 1、橙の破線と○: ランダムな $\vq$ と $\vk=\vq$ の 200 組の平均）は、差が大きいほど内積がおおむね小さくなる。赤と橙はほぼ重なる（差は最大で約 0.02）。互いに無関係なランダムな $\vq,\vk$（青: 200 組の平均、灰: 1 組）では、この傾向は見られない（\file{ch11/sweep_rope_decay.py}）。}
\label{fig:ma-rope}
\end{figure}

\figref{fig:ma-rope}(a) は式\eqref{eq:ma-rope-freq}の周期を描いたものです。
左端の組は約 6.3 トークンで一周し、右端の組は約 4.7 万トークンかけて一周します。
(b) は、クエリとキーの選び方を変えて、位置の差によって内積がどう変わるかを比べたものです（\file{ch11/sweep_rope_decay.py} で、4 通りの $\vq,\vk$ について計算しました）。
クエリとキーが同じ向きのとき（赤と橙の線）は、差が 0 のとき最大で、離れるほどおおむね小さくなっていきます。
赤と橙がほぼ重なるので、「全成分 1」という特別な選び方だけの性質ではありません。
一方、互いに無関係なランダムな $\vq,\vk$（青と灰の線）では、200 組の平均はどの差でもほぼ 0 で、1 組だけを見ても、差とともに下がるわけではありません。
つまり、「近いほど内積が大きい」のは、クエリとキーが似た向きのときに限った性質です。
RoPE の論文は、この性質を「離れたトークンほど結び付きが弱くなる」という直観に合うものとして説明しています\cite{su2024roformer}。
ただし、学習の前のランダムに近い $\vq,\vk$ に、近くを優先させるわけではありません。

\subsection*{実装}
\skimread{押さえる結論は「回転は、前もって作った cos と sin の表との積で計算できる。位置の表がないので、学習時より長い文脈にも計算としては使える（ただし、そのままでは性能が落ちる）」です。}

回転行列をそのまま掛けると、ほとんどが 0 の大きな行列積になって無駄です。
式\eqref{eq:ma-rot}を成分で書くと、組 $(x_1,x_2)$ は
$y_1=x_1\cos\phi-x_2\sin\phi$、$y_2=x_1\sin\phi+x_2\cos\phi$
になるので、これを成分ごとの積で計算します。

\snippet{RoPE の前計算と適用}{code:ma-rope}{ch11/modern_gpt.py}{rope}

3〜4 行目が式\eqref{eq:ma-rope-freq}で、\code{omega} の形は $(\dhead/2,)$ です。
5 行目で位置 $0,1,\dots$ の列と掛け合わせ、すべての位置と組の角度 $m\omega_i$ を $(T,\dhead/2)$ の表にします。
\code{[:, None]} で位置の列を形 $(T,1)$ にしてあるので、\code{omega} との積はブロードキャストでこの形になります。
この表は学習しないので、一度だけ計算して持っておきます。
11 行目の \code{x[..., 0::2]} は各組の 1 つ目の成分（式の $x_1$。0 から数えて偶数番目）、\code{x[..., 1::2]} は 2 つ目の成分 $x_2$ で、隣り合う 2 成分が 1 組になります。
12〜13 行目が上の回転の式そのものです。
14 行目で $(y_1,y_2)$ を交互に並べ直して元の形に戻します。

実際に性質を確かめましょう。

\snippet{RoPE の相対位置の性質を確かめる}{code:ma-rope-check}{ch11/rope.py}{rope-relative}
\outputfile{ch11/out/rope.txt}

実行結果の前半 5 行は、上の数値例を同じ関数で計算したものです（\file{ch11/rope.py} の前半）。
$m-n=+2$ の 2 行がどちらも $-0.4161\approx\cos2$ になっています。
$m=3,n=5$（差が $-2$）でも同じ値なのは、$\cos$ が偶関数だからです（2 次元で $\vq=\vk$ の特別な場合です）。
後半は 64 次元のランダムな $\vq,\vk$ です。
位置 (10, 7)、(110, 107)、(1000, 997) のスコアは、どれも差が 3 なので $-8.281$ でそろっています。
最後の桁の小さな違いは、角度 $m\omega_i$ が大きくなったときの浮動小数点の丸め誤差です。
差が 10 の 2 行も、$-7.482$ でそろっています。
最後の行の「回す前の q·k」$-11.434$ は、同じ位置どうし（差 0）のスコアと同じ値です。
両方を同じ角度だけ回しても、間の角度は変わらないからです。
差 0、3、10 でスコアが変わるので、位置の差がスコアに効いていることが分かります。
別の $\vq,\vk$ を使えば、別の値になります。
スコアには、中身と位置の差の両方が効いているのです。
なお、この組では離れたほうがスコアが大きくなっていますが、これは\figref{fig:ma-rope}(b) の青と灰の線と同じく、2 本の向きが無関係だからです。
また、回す前と回した後でノルムが 8.3705 のまま変わらないことも確かめられます。

RoPE には位置の表がないので、学習に使った長さより長い文脈にも\textbf{計算としては}適用できます。
ただし、そのままでは性能が落ちます。
回転は 1 周（$2\pi$）で元に戻るので、速い組は、学習した長さの中ですでに何周も回っています。
困るのはゆっくり回る組です。
たとえば 1 周に 628 トークンかかる組は、文脈長 128 の学習では 5 分の 1 周ほどの角度しか経験していません。
それより長い文脈では、学習で見たことのない角度とその組み合わせが現れるので、モデルの性能が落ちます。
そのため実際の LLM は、$\theta$ を 50 万（Llama 3）や 100 万（Qwen2.5、Mixtral）のように大きくして、ゆっくり回します。
学習の後半で長い文書を使い、文脈長を伸ばすこともします（\secref{sec:ma-beyond}）。

この節で分かったのは、RoPE がクエリとキーを位置に比例した角度だけ回すことで、注意のスコアに位置が「位置の差」の形でだけ入るようにする方法だということです（スコアは、クエリとキーの中身にもよります）。

\section{SwiGLU：ゲート付きの FFN}\label{sec:ma-swiglu}
\skimread{押さえる結論は「SwiGLU は、ゲートの枝が作る値を、中身の枝の各成分に掛ける倍率として使う FFN。行列は 3 枚で、中間次元を約 $\tfrac83d$ にすれば、パラメータ数は GPT-2 型の MLP とそろう」です。}

3 つ目の差分は、各位置のベクトルを位置ごとに変換する部品である FFN（第 10 章の MLP）です。
この節の問いは、「GELU の 2 層 MLP を、なぜ行列が 3 枚の形に置き換えたのか」です。

\subsection*{ゲートという考え方}

先に断っておくと、この置き換えは理論から導かれたものではありません。
比較実験で良い結果が出た形が、標準になりました（理由は下で述べます）。
それでも、何を変えたのかは式で正確に言えます。

第 10 章の MLP は $\vx\mapsto\mW_2^\T\gelu(\mW_1^\T\vx)$ でした。
中間の $4d$ 個の値は、それぞれ活性化関数を通って「どれだけ出力に効くか」が決まります。
このとき、各成分の効き方を決めているのは、その成分自身の値だけです。
そこで、効き方を\textbf{別の線形変換に決めさせる}ことを考えます。
2 つの線形変換 $\va=\mW_g^\T\vx$ と $\vb=\mW_u^\T\vx$ を作り、$\va$ を活性化関数に通した値を、$\vb$ の各成分に掛ける倍率として使うのです。
倍率が 0 に近ければ、対応する $\vb$ の成分はほとんど出力に届かず、大きければ大きく届きます。
この倍率を作る枝を\textbf{ゲート}（gate、門）と呼びます。
水道の蛇口にたとえると、ゲートが蛇口の開き具合、$\vb$ が流れてくる水に当たります。
ただし、このたとえには限界があります。
蛇口は閉（0 倍）から全開（1 倍）までですが、下で使う SiLU の値は 1 を超えることも、小さな負になることもあります。
このような層を GLU（gated linear unit）と呼びます。

Shazeer\cite{shazeer2020glu}は、GLU を Transformer の FFN に入れて活性化関数をいくつか比べ、SiLU（Swish とも呼ぶ）を使ったものが良い結果を出すことを示しました。
SiLU は $\mathrm{SiLU}(u)=u\,\sigmoid(u)$ で、$\sigmoid$ はシグモイド関数です。
形は GELU とよく似た、0 付近でなめらかに曲がる関数です。
たとえば $\mathrm{SiLU}(2)=2\times0.881=1.762$、$\mathrm{SiLU}(-2)=-2\times0.119=-0.238$ です。
SiLU の値は約 $-0.28$（$u\approx-1.28$ のとき）より小さくならず、上は限りなく大きくなります。
つまりゲートは、中身を止めるだけでなく、何倍かに強めたり、符号を変えて小さく通したりもします。

\begin{definition}[SwiGLU の FFN]
入力 $\vx\in\R^d$、中間次元を $\dff$ とし、$\mW_g,\mW_u\in\R^{d\times\dff}$、$\mW_d\in\R^{\dff\times d}$ とします。
\begin{equation}
\mathrm{SwiGLU}(\vx)=\mW_d^\T\bigl(\mathrm{SiLU}(\mW_g^\T\vx)\odot\mW_u^\T\vx\bigr).
\label{eq:ma-swiglu}
\end{equation}
$\mW_g$ はゲートの枝を作る行列（gate）、$\mW_u$ は中身の枝を作る行列で、どちらも次元を $d$ から $\dff$ に上げます（$\mW_u$ の u は up）。
$\mW_d$ は $\dff$ から $d$ に下げる行列です（d は down の頭文字で、次元の $d$ とは別）。
\end{definition}

\begin{example}[SwiGLU の小さな数値例]
$\dff=2$ で、ある $\vx$ から $\va=\mW_g^\T\vx=(2,\,-2)$、$\vb=\mW_u^\T\vx=(1,\,1)$ が作られたとします。
ゲートの値は $\mathrm{SiLU}(\va)=(1.762,\,-0.238)$ です。
これを $\vb$ と成分ごとに掛けると $(1.762\times1,\,-0.238\times1)=(1.762,\,-0.238)$ になります。
1 つ目の成分は約 1.8 倍に強められ、2 つ目は符号が変わって小さく通されました。
この 2 次元のベクトルを $\mW_d^\T$ で $d$ 次元に戻したものが、FFN の出力です。
\end{example}

\begin{figure}[tbp]
\centering
\begin{tikzpicture}[x=1mm,y=1mm,b/.style={box,font=\footnotesize,minimum width=17mm,minimum height=6mm}]
\node[b] (x) at (0,0) {$\vx$ \scriptsize$(d)$};
\node[b] (g) at (-16,13) {$\mW_g$ \scriptsize$d\!\to\!\dff$};
\node[b] (u) at (16,13) {$\mW_u$ \scriptsize$d\!\to\!\dff$};
\node[greenbox,font=\footnotesize,minimum width=14mm,minimum height=6mm] (s) at (-16,25) {SiLU};
\node[op] (m) at (0,34) {$\odot$};
\node[b] (d) at (0,45) {$\mW_d$ \scriptsize$\dff\!\to\! d$};
\draw[arr] (x) -- (g); \draw[arr] (x) -- (u);
\draw[arr] (g) -- (s); \draw[arr] (s) |- (m); \draw[arr] (u) |- (m);
\draw[arr] (m) -- (d); \draw[arr] (d) -- ++(0,9);
\node[lbl,anchor=east] at (-24,25) {ゲート（倍率）};
\node[lbl,anchor=west] at (1,52) {出力 \scriptsize$(d)$};
\node[lbl,anchor=west] at (17,25) {中身};
\end{tikzpicture}
\caption{SwiGLU の FFN。左の枝がゲートで、SiLU を通した値が、右の枝（中身）の各成分に掛ける倍率になる（倍率は 0〜1 に限らない）。$\odot$ は成分ごとの積。行列は 3 枚で、$\mW_g$ はゲートの枝を作る行列、$\mW_u$ は中身の枝を作る行列（どちらも $d$ 次元を中間次元 $\dff$ に上げる）、$\mW_d$ は $\dff$ から $d$ に下げる行列。}
\label{fig:ma-swiglu}
\end{figure}

\subsection*{パラメータ数をそろえる}

行列が 2 枚から 3 枚に増えたので、同じ $\dff=4d$ のままではパラメータが 1.5 倍になります。
GPT-2 型の MLP は $2\cdot d\cdot4d=8d^2$、SwiGLU は $3d\,\dff$ です。
両者を等しくするには
\[
3d\,\dff=8d^2 \quad\Longrightarrow\quad \dff=\tfrac{8}{3}d\approx2.67d
\]
とすればよいことが分かります。
LLaMA はこの $\tfrac{8}{3}d$ を使いました\cite{touvron2023llama}。
実際の設定値は、計算の効率のために切りのよい数に丸められます。
Llama 2 7B（7B はパラメータが約 70 億という意味。B は billion＝10 億）は $d=4096$ に対して $\dff=11008$（$11008/4096\approx2.69$ で $2.69d$）、Llama 3 8B は $14336$（$3.5d$）、Qwen3-8B は $12288$（$3d$）です。
Llama 2 は $\tfrac83d$ に近く、Llama 3 と Qwen3 は、そろえる値より大きめに取っています。
本章の現代版ミニGPT では、$d=128$ に対して $\tfrac83\times128\approx341$ を 8 の倍数に丸めた 344 を使います。


\figref{fig:ma-swiglu}の左の枝が「倍率」、右の枝が「中身」です。
2 つの枝が同じ $\vx$ から作られているので、モデルは「入力のこの特徴があるときだけ、この中身を通す」という掛け算の関係を、1 層で表せます。
なぜこれがよく効くのかについて、Shazeer 自身は論文で「説明はできない」と率直に書いています\cite{shazeer2020glu}。
SwiGLU は理論からではなく、比較実験の結果として標準になった部品です。

\snippet{SwiGLU}{code:ma-swiglu}{ch11/modern_gpt.py}{swiglu}

4〜6 行目が 3 枚の行列で、すべてバイアスなしです。
9 行目が式\eqref{eq:ma-swiglu}そのもので、\code{F.silu(self.w_gate(x))} がゲート、\code{self.w_up(x)} が中身、\code{*} が成分ごとの積 $\odot$ です。
Gemma 系は SiLU の代わりに GELU を使う GeGLU を採用しています。

この節で分かったのは、SwiGLU が「ゲートの枝が作る倍率を、中身の枝に成分ごとに掛ける」FFN で、中間次元を約 $\tfrac83 d$ にすればパラメータ数を GPT-2 型とそろえられることです。

\section{バイアスと重み共有}\label{sec:ma-bias}
\skimread{押さえる結論は「現代のモデルは、線形層と正規化のバイアスをほぼ一律に取り除く。埋め込みと出力の重み共有は、小さいモデルでは使い、大きいモデルでは使わない」です。}

4 つ目の差分は小さな整理です。
この節の問いは、「線形層のバイアスと、埋め込みと出力の重み共有は、現代のモデルでどう扱われているか」です。

\textbf{バイアスなし}。
GPT-2 の線形層と LayerNorm には、バイアス（ずらし）の項がありました。
現代の多くのモデルは、これをすべて取り除いています。
理由は主に 2 つです。
一つは、取り除いても性能がほとんど変わらないことが、経験的に知られていることです。nanoGPT の実装にも、バイアスなしのほうが「少し良く、少し速い」という注記があります\cite{karpathy2022nanogpt}。
もう一つは、部品が減るぶん、モデルが単純になり、計算もわずかに軽くなることです。
実は本書の \code{llmbook} のミニGPT も、線形層はすでにバイアスなしで作ってあります（\chapref{ch:transformer}）。
Qwen3 は前の世代にあったクエリ・キー・バリューのバイアスを外しました\cite{yang2025qwen3}。
例外もあり、gpt-oss は注意の線形層にバイアスを持っています。

\textbf{重み共有}。
第 10 章では、出力ヘッドがトークン埋め込みの行列 $\mE$ と同じ行列を使う重み共有（weight tying）\cite{press2017using}を説明しました。
現代では、これはモデルの大きさで分かれます。
現在の LLM の語彙は 15 万〜26 万と大きいので、小さなモデルでは、$Vd$ の埋め込みが全体のかなりの割合を占めます。
共有すれば、その分を節約できます。
Gemma や、Qwen3 の 4B 以下のモデルは共有しています。
一方、Llama 3 や Qwen3 の 8B 以上、DeepSeek-V3 は共有していません。
モデルが大きいと埋め込みの割合が小さくなり、入力側と出力側で別の行列を持つ自由度のほうが役立つと判断されているのです。
本章の比較実験では、第 10 章のモデルと条件をそろえるため、共有したまま比べます。

この節で分かったのは、バイアスはほぼ一律に取り除かれ、重み共有はモデルの大きさに応じて選ばれていることです。

\begin{quiz}
\tabref{tab:ma-diff}の上の 4 行（RMSNorm、RoPE、SwiGLU、バイアスなし）は、それぞれモデルのパラメータ数を増やすでしょうか、減らすでしょうか、それともほぼ変えないでしょうか。
\end{quiz}

答え合わせをしましょう。
RoPE は位置の表 $\mP$（$T_{\max}d$ 個）をなくすので、パラメータは減ります。
RMSNorm は $\bm\beta$ がなくなる分、バイアスなしはバイアスの分だけ減ります。
SwiGLU は、$\dff=\tfrac83 d$ にそろえれば\textbf{ほぼ変わりません}。
つまり、減るものはあっても、増えるとしても丸めによるごくわずかな分だけで（\secref{sec:ma-assemble}）、どれもパラメータ数を大きくは変えません。
現代の設計の改良の多くは、同じパラメータ数で「より学習しやすく、より長い文脈を扱える」形を探したものなのです。

\subsection*{ここまでの道のり}

ここで一度、現在地を確かめておきます。
\begin{itemize}
\item 正規化は RMSNorm に、位置の情報は RoPE に、FFN は SwiGLU に置き換えました。
\item どれもパラメータ数はほとんど変えずに、学習のしやすさや長い文脈への対応を改善する変更でした。
\item ここからは\textbf{生成を速くする}工夫（KV キャッシュと GQA）と、\textbf{パラメータを安く増やす}工夫（MoE）に進みます。
\end{itemize}

\section{KV キャッシュ：生成の無駄をなくす}\label{sec:ma-kvcache}

ここからは、学習ではなく\textbf{生成}の話です。
第 10 章の \code{generate} は、1 トークンを生成するたびに、それまでの系列全体をモデルに入れ直していました。
この節の問いは、「その計算のうち、どこが無駄で、どうすれば省けるのか」です。

\subsection*{何が無駄なのか}

プロンプトが「吾輩は」で、モデルがその続きを「猫である。名前は…」と生成していく場面を、1 ステップずつ追ってみましょう。
本書の BPE（\chapref{ch:tokenizer}）では、「吾輩は」は 1 個のトークンです。
続きは「猫」「である。」「名」「前」「はまだ」の順のトークンになります。
この続きは説明のために選んだ列です（\secref{sec:tf-walkthrough}で基盤モデルに「＼n　吾輩は」を与え、貪欲法で書かせた実際の続きは、「主人」から始まる列でした）。
列が違うので、この節の「位置 1」は「吾輩は」で、\secref{sec:tf-walkthrough}の「位置 1」（＼n　）とは別のトークンです。
省ける計算の数え方は、どの列でも同じです。
1 ステップでは、第 10 章で見たように、最後の位置のロジットから確率を作って次のトークンを 1 個選び、列の末尾につなぎます。
第 10 章の \code{generate} は、各ステップで列の全体をモデルに入れ直していました。
1 回目は「吾輩は」の 1 トークンを入れて、2 個目のトークン「猫」を得ます。
2 回目は「吾輩は」「猫」の 2 トークンを入れて、3 個目の「である。」を得ます。
3 回目は 3 トークンを入れて……と続きます。

ここで、2 回目の計算の中身を考えます。
因果マスク（\chapref{ch:attention}）があるので、位置 1 の計算は、位置 2 のトークンが何であっても変わりません。
各層の位置 1 のキーとバリューは、1 回目に計算したものとまったく同じなのです。
3 回目なら位置 1 と 2 が、4 回目なら位置 1〜3 が、前の回と同じです。
それなのに、毎回すべての位置について計算し直しています。
一般に、プロンプトの長さを $T_0$ として $n$ 個のトークンを生成すると、モデルに通すトークンの数は、合わせて $T_0+(T_0+1)+\dots+(T_0+n-1)=nT_0+n(n-1)/2$ になります。
$n$ 個の項のそれぞれに $T_0$ があり、残りの $0+1+\dots+(n-1)$ が $n(n-1)/2$ だからです。
「吾輩は」の例（$T_0=1$、$n=5$）なら $5+10=15$ 個です。
第 2 項があるので、長く生成すると、この数はほぼ $n$ の 2 乗に比例して増えていきます。

では、新しい位置の計算に、過去のどの量が必要でしょうか。
位置 $t$ の注意の出力は、クエリ $\vq_t$ と、自分を含むそこまでのすべての位置のキー $\vk_1,\dots,\vk_t$、バリュー $\vv_1,\dots,\vv_t$ から計算されます。
ほかの位置のクエリは使いません。
FFN と正規化は位置ごとに別々に計算され、ほかの位置の値を使いません（位置の間で情報をやり取りするのは注意だけでした）。
そのため、2 層目以降が過去の位置について必要とするのも、その層のキーとバリューだけです。
過去の位置の値は、因果マスクのおかげで後のトークンに左右されません。
各層に取っておいたキーとバリューは、その層より下の計算の結果をすでに含んだものとして、そのまま使えます。
したがって、\textbf{各層の過去のキーとバリューだけを取っておけば}、新しいトークンについての 1 個分の順伝播で済みます。
ただし注意の中では、新しいクエリ 1 個と、取っておいた $t$ 個のキーとの内積、$t$ 個のバリューの重み付き和を、毎回計算します。
この部分だけは、文脈が長いほど少しずつ重くなります。
省けるのは、過去の位置の埋め込み・射影・FFN と、過去の位置のクエリの注意です。

\begin{definition}[KV キャッシュ]
\term{KV キャッシュ}{key--value cache}とは、自己回帰生成の間、各層・各位置で計算したキーとバリューを保存しておき、
次のステップでは新しいトークンの分だけキーとバリューを計算して保存済みのものに追加し、注意を計算する仕組みです。
\end{definition}

\begin{figure}[tbp]
\centering
\begin{tikzpicture}[x=1mm,y=1mm,
  kc/.style={cell,minimum width=8.6mm,minimum height=4.4mm,font=\tiny},
  qc/.style={cell,minimum width=6mm,minimum height=3.8mm,font=\tiny}]
\node[font=\scriptsize\bfseries] at (31,11) {(a) キャッシュなし};
\node[font=\tiny] at (31,7.8) {毎回、全部のトークンを入れ直す};
\node[font=\scriptsize\bfseries] at (83.5,11) {(b) キャッシュあり};
\node[font=\tiny] at (83.5,7.8) {新しい 1 トークンだけを入れる};
\node[font=\scriptsize\bfseries,align=center] at (118,9.5) {選んだ\\トークン\\[-0.3ex]{\tiny\mdseries (a)(b) 共通}};
\draw[black!20] (59.25,9) -- (59.25,-77);
\foreach \s/\out in {1/猫,2/である。,3/名,4/前,5/はまだ} {
  \pgfmathsetmacro{\y}{-(\s-1)*16.5}
  \node[font=\scriptsize,anchor=east] at (7.5,\y-1.5) {ステップ \s};
  \pgfmathsetmacro{\xs}{12+(\s-1)*9.5}
  \pgfmathsetmacro{\xsb}{64.5+(\s-1)*9.5}
  \foreach \j/\tk in {1/吾輩は,2/猫,3/である。,4/名,5/前} {
    \ifnum\j>\s\relax\else
      \pgfmathsetmacro{\xa}{12+(\j-1)*9.5}
      \pgfmathsetmacro{\xb}{64.5+(\j-1)*9.5}
      \draw[black!60,thin] (\xs,\y-5.6) to[out=90,in=-90,looseness=0.5] (\xa,\y-2.2);
      \draw[black!60,thin] (\xsb,\y-5.6) to[out=90,in=-90,looseness=0.5] (\xb,\y-2.2);
      \node[font=\tiny] at (\xa,\y+4.6) {\tk};
      \node[kc,fill=figorange] at (\xa,\y) {$\vk_{\j},\vv_{\j}$};
      \ifnum\j=\s
        \node[qc,fill=figorange] at (\xa,\y-7.5) {$\vq_{\j}$};
        \node[font=\tiny\bfseries] at (\xb,\y+4.6) {\tk};
        \node[kc,fill=figorange] at (\xb,\y) {$\vk_{\j},\vv_{\j}$};
        \node[qc,fill=figorange] at (\xb,\y-7.5) {$\vq_{\j}$};
      \else
        \node[qc,fill=black!8,draw=black!25,text=black!40] at (\xa,\y-7.5) {$\vq_{\j}$};
        \node[font=\tiny,text=black!35] at (\xb,\y+4.6) {\tk};
        \node[kc,fill=figblue] at (\xb,\y) {$\vk_{\j},\vv_{\j}$};
      \fi
    \fi
  }
  \node[font=\scriptsize,anchor=west] at (109,\y-1.5) {$\to$ \out};
}
\node[font=\tiny,align=center] at (31,-78.5) {計算した $\vk,\vv$：$1+2+3+4+5=15$ 組};
\node[font=\tiny,align=center] at (83.5,-78.5) {計算した $\vk,\vv$：5 組（読んだのは $0+1+2+3+4=10$ 組）};
\begin{scope}[shift={(4,-85)}]
\node[cell,minimum width=4mm,minimum height=3mm,fill=figorange] at (0,0) {}; \node[font=\tiny,anchor=west] at (2.5,0) {そのステップで計算};
\node[cell,minimum width=4mm,minimum height=3mm,fill=figblue] at (25,0) {}; \node[font=\tiny,anchor=west] at (27.5,0) {キャッシュから読む};
\node[cell,minimum width=4mm,minimum height=3mm,fill=black!8,draw=black!25] at (50,0) {}; \node[font=\tiny,anchor=west] at (52.5,0) {前のステップと同じ値を計算し直す};
\draw[black!60,thin] (81,-1.5) -- (83,1.5); \node[font=\tiny,anchor=west] at (84,0) {新しいクエリが内積をとるキー};
\end{scope}
\end{tikzpicture}
\caption{「吾輩は」から「猫である。名前は…」を生成する計算を、KV キャッシュなし (a) とあり (b) で 1 ステップずつ比べたコマ送りの図（模式図。1 つの層だけを描いた）。各コマの上の字は、そのステップでモデルに入れるトークン（(b) の薄い字は入れない）。中段のマスはその位置のキーとバリュー、下段はその位置のクエリ。右端は、最後の位置のロジットから選んだトークンで、(a) と (b) で同じになる。トークンの区切りは本書の BPE（\chapref{ch:tokenizer}）による。この列は、\chapref{ch:transformer}で実測した「吾輩は」からの生成（\figref{fig:tf-wagahai}、\figref{fig:tf-wagahai-steps}）とは別の、説明のための例で、本書の小さなモデルがこう生成するとは限らない。}
\label{fig:ma-kvcache}
\end{figure}

\figref{fig:ma-kvcache}の各コマを、左右で見比べてください。
(a) では、ステップが進むたびに、入れるトークンも、橙のキーとバリューのマスも 1 つずつ増えます。
灰色のクエリは、前のステップと同じ注意を、もう一度計算します。
最後の層では、その結果を次のトークンを選ぶのに使いません。
次のトークンを選ぶのに使うのは、最後の位置の出力だけだからです。
それより下の層では、結果を次の層のその位置のキーとバリューを作るのに使いますが、できる値は前のステップと同じです。
(b) では、2 回目以降に入れるのは新しい 1 トークンだけで、橙のキーとバリューのマスは各コマに 1 組しかありません。
過去の位置のキーとバリュー（青）は、メモリから読むだけです。
細い線は、新しいクエリが内積をとるキーです。
線の本数と行き先は (a) と (b) でまったく同じなので、注意の出力も、選ばれるトークンも同じになります。
橙のマスを数えると、(a) は $1+2+3+4+5=15$ 組、(b) は 5 組です。
その代わり、(b) では青と橙のマス、つまりキャッシュに取っておくキーとバリューが、ステップごとに 1 トークン分ずつ増えていきます。
KV キャッシュは、\textbf{計算を減らす代わりにメモリを使う}方法です。
このように、計算量とメモリの一方を減らすと他方が増える関係を、\term{トレードオフ}{trade-off}と呼びます。
1 トークン分のキャッシュが何バイトになるかは、\secref{sec:ma-gqa}で実際のモデルの数を入れて計算します。

\subsection*{実装}

キャッシュの箱は簡単です。
層ごとに 1 つ用意し、新しいキーとバリューを位置の軸（系列の長さの軸）につないでいくだけです。

\snippet{KV キャッシュの箱}{code:ma-kvcache}{ch11/modern_gpt.py}{kvcache}

箱は最初は空（4 行目の \code{None}）で、1 回目の \code{update} では渡されたものをそのまま入れます（7〜8 行目）。
10〜11 行目の \code{torch.cat(..., dim=2)} が、形 $(B,H_{\mathrm{kv}},t,\dhead)$ の位置の軸（左から 3 番目。0 から数えて 2 なので \code{dim=2}）にキーとバリューを追加する部分です。
$t$ は、そこまでにたまった位置の数です（コードのコメントの「時刻の軸」は、この位置の軸のことです）。
$H_{\mathrm{kv}}$ はキーとバリューのヘッドの数で、次の節で説明します（今はヘッド数 $H$ と同じだと思って構いません）。

これを使うのが注意の層（次の節の\coderef{code:ma-gqa}の 16〜17 行目）です。
新しいトークンのキーとバリューを計算した直後に、過去の分とつなぎ、つないだもの全体に対してクエリで注意をとります。
RoPE はキャッシュに入れる\textbf{前}にかけてあるので、過去のキーを回し直す必要はありません。
ここで、モデルに「いま何番目の位置か」を教える必要があります。
新しく入れるのは 1 トークンだけなので、そのままでは位置 0 だと誤解されてしまうからです。
コードでは位置を 0 から数えます。
\figref{fig:ma-kvcache}と本文の位置 1, 2, \dots は、コードの位置 0, 1, \dots に当たります。
それが次の \code{forward} の引数 \code{start} です。

\snippet{現代版ミニGPT の順伝播}{code:ma-forward}{ch11/modern_gpt.py}{forward}

これは、本章の部品で組んだ現代版ミニGPT（\code{ModernGPT}）の順伝播です（組み立ては\secref{sec:ma-assemble}）。
\code{self.cos} と \code{self.sin} は、RoPE で回す角度の cos と sin を前もって計算した表で、1 行が 1 つの位置に対応します（\coderef{code:ma-rope}）。
3〜4 行目で、この表から位置 \code{start} から \code{start+T-1} の行だけを取り出しています。
キャッシュを使う生成では \code{start} に「これまでのトークン数」を渡すので、新しいトークンは正しい位置の角度で回されます。
5 行目はトークン埋め込みです。
第 10 章の \code{GPT} と違って、位置埋め込みを足していません。
位置の情報は、各ブロックの注意の中で RoPE が与えるからです。
6〜7 行目で、各ブロックに角度の表と、その層のキャッシュ \code{caches[i]} を渡します。
8 行目の最後の正規化と出力ヘッド、9〜12 行目の損失は、第 10 章の \code{GPT} と同じ形です。
位置の表 $\mP$ を足す GPT-2 型でも、位置 \code{start} からの行を表から取り出して足せば、同じように KV キャッシュを使えます。

生成の関数は次のとおりです。

\snippet{キャッシュあり・なしの貪欲生成}{code:ma-generate}{ch11/modern_gpt.py}{generate}

1 行目の \code{@torch.no_grad()} は、第 10 章の \code{generate} と同じく、勾配のための記録をとらない指定です。
4 行目で、層の数だけ空の箱を用意します。
5 行目で、まずプロンプト全体を一度に処理します。
このときキャッシュには、プロンプトの全位置のキーとバリューがたまります。
7 行目は、ロジットが最大のトークンを選ぶ貪欲法（\chapref{ch:lm}）で、8 行目でそれを系列 \code{idx} の末尾につなぎます。
その後、キャッシュありなら 10 行目で新しい 1 トークン（\code{nxt}、形 $(B,1)$）だけを入れ、なしなら 12 行目で系列全体を入れ直します。
10 行目の \code{start=idx.shape[1] - 1} は、新しいトークンの位置です。
\code{idx} にはもう新しいトークンをつないであるので、長さから 1 を引くと、0 から数えた位置になります。
プロンプトが 16 トークン（位置 0〜15）なら、最初の新しいトークンは位置 16 で、そのとき \code{idx} の長さは 17 です。
なお、このループは最後の回にも、使わないロジットのための順伝播を 1 回します（返すのは \code{idx} だけです）。
そのため実際に通すトークンの数は、上の式より 1 回分だけ多くなります。

\subsection*{正しさと速さを測る}

キャッシュは「計算のやり方」を変えるだけなので、結果は同じになるはずです。
まずそれを確かめ、それから速さを比べます。
モデルは学習前の現代版ミニGPT（4 層、$d=192$、4 ヘッド。キー・バリューのヘッドも 4 個で、GQA はまだ使っていません）で、プロンプトは 16 トークンです（$T_0=16$）。
速さを測るだけなら、学習済みである必要はありません。
\file{ch11/kv_cache.py} の前半では、キャッシュありとなしで 64 トークンを生成して列を比べ、さらに 1 トークンずつ追加して計算したロジットと、全体を一度に計算したロジットを比べます。
後半が次の時間の測定です。

\snippet{キャッシュの有無で時間を比べる}{code:ma-kv-bench}{ch11/kv_cache.py}{kv-bench}
\outputfile{ch11/out/kv_cache.txt}

1 行目と 2 行目が正しさの確認です。
生成したトークン列は完全に一致し、1 トークンずつ追加して計算したロジットと、全体を一度に計算したロジットの差は最大でも $8\times10^{-7}$ でした。
これは浮動小数点の足し算の順序が違うことによる丸め誤差の大きさで、計算としては同じものだと言えます。

3 行目以降が速さです（各 3 回の中央値）。
\figref{fig:ma-kv-speed}にグラフにしました。
キャッシュありの時間は、新しいトークンの数にほぼ比例しています（32 個で 74 ms、480 個で 1,171 ms。1 トークンあたり約 2.3〜2.4 ms）。
1 トークンあたりの時間が 2.31 ms から 2.44 ms へ少しずつ増えるのは、新しいクエリが内積をとるキーの数が、文脈とともに増えるからです。
キャッシュなしの時間は、それより速く増えています。
新しいトークンの数が 15 倍（32 → 480）になると、時間は約 36 倍（267 ms → 9,553 ms）になりました。
比例（15 倍）より速く増えるのは、キャッシュなしでモデルに通すトークンの数が、$n$ の 2 乗の項を持つからです。
ただし 36 倍は、数から予想される倍率（\tabref{tab:ma-kv-count}で約 122 倍）よりずっと小さな値です。
その理由は、表の後で説明します。
その結果、速度の比は 32 個の 3.6 倍から 480 個の 8.2 倍へと、長く生成するほど大きくなります。

\begin{table}[tbp]
\centering
\caption{モデルに通すトークンの数（本文の式で数えた、必要な最小の個数）と、測った時間。$T_0=16$。キャッシュなしは $nT_0+n(n-1)/2$ 個、ありは $T_0+(n-1)$ 個。数の比は 21〜248 倍あるのに、時間の比は 3.6〜8.2 倍にとどまる。\coderef{code:ma-generate}は最後に結果を使わない順伝播を 1 回するので、実際に通す数はわずかに多い（$n=32$ で 1,056 個と 48 個）。時間はその分も含めて測った。}
\label{tab:ma-kv-count}
\small
\begin{tabular}{@{}rrrrrrr@{}}
\toprule
 & \multicolumn{3}{c}{通すトークンの数} & \multicolumn{3}{c}{時間（ms）} \\
\cmidrule(lr){2-4}\cmidrule(l){5-7}
新しいトークン $n$ & なし & あり & 比 & なし & あり & 比 \\
\midrule
32 & 1,008 & 47 & 21 & 267 & 74 & 3.6 \\
64 & 3,040 & 79 & 38 & 564 & 147 & 3.8 \\
128 & 10,176 & 143 & 71 & 1,328 & 290 & 4.6 \\
256 & 36,736 & 271 & 136 & 3,592 & 606 & 5.9 \\
480 & 122,640 & 495 & 248 & 9,553 & 1,171 & 8.2 \\
\bottomrule
\end{tabular}
\end{table}

\tabref{tab:ma-kv-count}の左半分は、通すトークンの数を先ほどの式で数えたものです。
キャッシュなしは、$n=32$ なら $32\times16+32\times31/2=512+496=1{,}008$ 個、$n=480$ なら $7{,}680+114{,}960=122{,}640$ 個です。
キャッシュありは、プロンプトの 16 個の後は 1 個ずつなので、$16+(n-1)$ 個で済みます。
ところが、右半分の時間は、数ほどには違いません。
$n=32$ では、数は 21 倍違うのに、時間は 3.6 倍「しか」違いません。
$n$ を 32 から 480 にしたときも、キャッシュなしの数は約 122 倍になるのに、時間は約 36 倍です。
これは、本書のモデルが小さいためです。
小さなモデルでは、30 トークンを処理するのも 1 トークンを処理するのも、Python や PyTorch の呼び出しの手間が時間の大半を占め、計算量の差がそのまま時間に出ません。
大きなモデルで長い文章を生成する場合は、計算量の差がそのまま時間の差になるので、キャッシュなしの生成は現実的ではありません。
実際の LLM の推論は、ほぼ例外なく KV キャッシュを使っています。
なお、CPU は他の作業と共有しているので、測るたびに数割ぶれることがあります。

\begin{figure}[tbp]
\centering
\begin{tikzpicture}
\begin{axis}[xlabel={生成した新しいトークンの数},ylabel={時間（ms）},legend pos=north west,
  scaled y ticks=false,ytick={0,2000,4000,6000,8000,10000},
  yticklabel style={/pgf/number format/.cd,1000 sep={,}}]
\addplot[bookgray,mark=*,mark size=1.3pt] table[col sep=comma,x=n,y=nocache]{code/ch11/out/kv_speed.csv};
\addlegendentry{キャッシュなし}
\addplot[bookaccent,mark=square*,mark size=1.3pt] table[col sep=comma,x=n,y=cache]{code/ch11/out/kv_speed.csv};
\addlegendentry{キャッシュあり}
\end{axis}
\end{tikzpicture}
\caption{KV キャッシュの有無による生成時間（4 層・$d=192$ のモデル、プロンプト 16 トークン、CPU 2 スレッド、3 回の中央値）。キャッシュありはほぼ直線、キャッシュなしは長くなるほど急に増える。}
\label{fig:ma-kv-speed}
\end{figure}

\subsection*{キャッシュが効くのはどんなときか}

ここまでの結果を、「どんなときにキャッシュが効くのか」という問いでまとめ直します。
\figref{fig:ma-kv-when}(a) は、1 ステップごとにモデルに通すトークンの数を、式で数えたものです（$T_0=16$）。
キャッシュなしでは、$s$ ステップ目に $T_0+s-1$ 個を通します。
キャッシュありでは、最初のステップの $T_0$ 個の後は、1 個ずつです。
2 本の線の差が、そのステップでキャッシュが省く計算の量です。
差はステップごとに広がるので、\textbf{生成が長いほど}キャッシュはよく効きます。
これが 1 つ目の場面です。

2 つ目は、\textbf{プロンプトが長いとき}、とくに\textbf{同じ前置きを何度も使うとき}です。
\file{ch11/sweep_kv_prompt.py} で、前置き（ここでは、生成の前に与えるプロンプトの全体）の長さを 16〜448 トークンの 5 通りに変え、新しいトークンを 64 個生成する時間を 3 つのやり方で測りました（\figref{fig:ma-kv-when}(b)）。
キャッシュなし（灰）は、64 回のステップのたびに前置きを計算し直すので、前置きが長いほど時間が大きく増えます。
キャッシュあり（赤）は、前置きを最初に 1 回だけ計算するので、増え方がずっと小さくなります。
「使い回し」（青）は、前置きのキャッシュを前もって計算して取っておき、生成の部分だけを測ったものです。
チャットのシステムプロンプトのように、同じ前置きで何度も生成するなら、前置きの計算そのものを毎回せずに済みます。
（\chapref{ch:eval-serving}では、前置きという語を、取っておいた列と先頭から一致する部分の意味で使います。この実験では、プロンプトの全体が一致しています。）
この使い回しをプレフィックスキャッシュと呼び、\chapref{ch:eval-serving}で詳しく扱います。
本書の小さなモデルでは、前置きの計算は 1 回の順伝播で済むので、赤と青の差は小さめで、測定のぶれと同じ程度の点もあります（赤の 256 と 448 がほぼ同じ時間になっているのも、ぶれによるものです）。

反対に、プロンプトも生成も短いときは、キャッシュの効果は小さくなります（\tabref{tab:ma-kv-count}の $n=32$ では 3.6 倍）。
そして、どの場面でもキャッシュはメモリを使います。
その大きさを、次の節で数えます。

\begin{figure}[tbp]
\centering
\begin{tikzpicture}
\begin{axis}[width=0.49\linewidth,height=0.42\linewidth,
  xlabel={ステップ $s$},ylabel={通すトークンの数},xmin=0,xmax=480,ymode=log,ymin=0.7,ymax=1500,
  legend pos=north west,legend style={font=\tiny},legend cell align=left,
  title={\footnotesize (a) 1 ステップに通す数（式、$T_0=16$）},title style={yshift=-1mm}]
\addplot[bookgray,thick,domain=1:480,samples=120] {16+x-1};
\addlegendentry{なし：$T_0+s-1$}
\addplot[bookaccent,thick,mark=*,mark size=1pt,mark indices={1}] coordinates {(1,16) (2,1) (480,1)};
\addlegendentry{あり：最初だけ $T_0$}
\end{axis}
\end{tikzpicture}\hfill
\begin{tikzpicture}
\begin{axis}[width=0.49\linewidth,height=0.42\linewidth,
  xlabel={前置きの長さ $T_0$（トークン）},ylabel={64 個の生成の時間（ms）},ymode=log,ymin=150,ymax=40000,
  log ticks with fixed point,yticklabel style={/pgf/number format/.cd,1000 sep={,}},
  legend pos=north west,legend style={font=\tiny},legend cell align=left,
  title={\footnotesize (b) 前置きの長さを変える（実測）},title style={yshift=-1mm}]
\addplot[bookgray,mark=*,mark size=1.2pt] table[col sep=comma,x=T0,y=nocache]{code/ch11/out/sweep_kv_prompt.csv};
\addlegendentry{なし}
\addplot[bookaccent,mark=square*,mark size=1.2pt] table[col sep=comma,x=T0,y=cache]{code/ch11/out/sweep_kv_prompt.csv};
\addlegendentry{あり}
\addplot[bookblue,mark=triangle*,mark size=1.5pt] table[col sep=comma,x=T0,y=reuse]{code/ch11/out/sweep_kv_prompt.csv};
\addlegendentry{前置きを使い回す}
\end{axis}
\end{tikzpicture}
\caption{KV キャッシュが効く場面。(a) 1 ステップごとにモデルに通すトークンの数を式で数えた値。キャッシュなしは $s$ とともに増え続け、ありは 2 ステップ目から 1 個になる。(a)(b) とも縦軸は対数。(b) 前置きの長さを 5 通りに変え、新しいトークンを 64 個生成した時間（\file{ch11/sweep_kv_prompt.py}。モデルは\tabref{tab:ma-kv-count}と同じ、CPU 2 スレッド、5 回の中央値）。「使い回す」は、前置きのキャッシュを前もって作っておき、生成の部分だけを測った値。\tabref{tab:ma-kv-count}とは別の時に測ったので、時間の絶対値は表と直接は比べられない（計算機を共有しているため）。同じ条件の\tabref{tab:ma-kv-count}の $n=64$ の行は、なし 564 ms・あり 147 ms で、この図の $T_0=16$ の点（2,183 ms・275 ms）とは約 2〜4 倍違い、なし／ありの比も 3.8 倍と 7.9 倍で違う。比べてよいのは、この図の中の線どうしだけである。}
\label{fig:ma-kv-when}
\end{figure}

この節で分かったのは、KV キャッシュが「過去のキーとバリューは変わらない」ことを利用して、生成の順伝播を新しいトークンの分だけにする仕組みであり、結果を変えずに速さを大きく改善すること、その代わりに生成が進むほどメモリを使うことです。

\section{GQA：キャッシュを小さくする}\label{sec:ma-gqa}

KV キャッシュで計算は減りました。
代わりに、キャッシュがメモリを使います。
この節の問いは、「キャッシュはどれくらい大きくなるのか、そしてどうすれば小さくできるのか」です。

\subsection*{キャッシュの大きさを数える}

1 トークンあたり、各層で、ヘッドごとにキーとバリューのベクトル（それぞれ $\dhead$ 次元）を 1 つずつ保存します（\figref{fig:ma-kvcache}の 1 マスが、1 つの層の 1 トークン分です）。
キーとバリューのヘッドが $H_{\mathrm{kv}}$ 個、1 ヘッドの次元が $\dhead$、層が $L$ 個、1 つの数を $b$ バイトで保存するとすると、
\begin{equation}
\text{1 トークンあたりの KV キャッシュ}=2\times L\times H_{\mathrm{kv}}\times\dhead\times b\ \text{バイト}
\label{eq:ma-kv-mem}
\end{equation}
です。
先頭の 2 は、キーとバリューの 2 つ分です。
これに文脈のトークン数と、同時に処理する系列の数（バッチサイズ）を掛けたものが、キャッシュ全体の大きさになります。

数値を入れてみましょう。
Llama 2 7B は $L=32$、$H_{\mathrm{kv}}=32$、$\dhead=128$ で、BF16（1 つの数が 2 バイト）で保存すると
$2\times32\times32\times128\times2=524{,}288$ バイト、つまり 1 トークンあたり 512 KiB です。
BF16 は 1 つの数を 16 ビットで表す形式で、通常の 32 ビット（4 バイト）の半分の大きさです（詳しくは\chapref{ch:gpu}）。
KiB は $2^{10}=1024$ バイト、GiB は $2^{30}$ バイトで、GB（$10^9$ バイト）とは約 7\% 違います。
文脈 8,192 トークン（Llama 2 自身の上限は 4,096 ですが、後で Llama 3 と比べるための長さ）なら 4 GiB になります。
モデルの重み（約 67 億パラメータ × 2 バイト ≈ 13 GB）の 3 分の 1 近くを、たった 1 本の会話のキャッシュが使うのです。
同時に 10 人の利用者に応えれば 40 GiB です。
GPU のメモリは限られているので（\chapref{ch:gpu}）、キャッシュの大きさが、同時に処理できる利用者の数と文脈の長さを直接制限します。

\subsection*{キーとバリューのヘッドを共有する}

式\eqref{eq:ma-kv-mem}のうち、設計で減らせるのは $H_{\mathrm{kv}}$ です。
$L$ や $\dhead$ を減らすと、注意も FFN も含めたモデルの全体が小さくなり、$b$ は数の形式で決まってしまうからです。
$H_{\mathrm{kv}}$ を減らして小さくなるのは、キーとバリューの射影だけで、クエリの側は変わりません。
マルチヘッド注意（\chapref{ch:attention}）では、クエリ・キー・バリューのヘッドの数はどれも $H$ でした。
では、キーとバリューのヘッドを減らし、複数のクエリのヘッドで共有したらどうでしょうか。

極端な例が、キーとバリューのヘッドを 1 つにする \term{MQA}{multi-query attention}\cite{shazeer2019fast}です。
キャッシュは $1/H$ になりますが、表現力が落ちて性能が下がることがありました。
その中間が \term{GQA}{grouped-query attention}\cite{ainslie2023gqa}です。
クエリのヘッドを $H_{\mathrm{kv}}$ 個のグループに分け、同じグループのクエリのヘッドは同じキー・バリューのヘッドを使います。
$H_{\mathrm{kv}}=H$ なら普通のマルチヘッド注意（MHA）、$H_{\mathrm{kv}}=1$ なら MQA です。

\begin{figure}[tbp]
\centering
\begin{tikzpicture}[x=1mm,y=1mm,
  q/.style={cell,minimum width=3.6mm,minimum height=4.5mm,fill=figblue},
  kv/.style={cell,minimum width=3.6mm,minimum height=4.5mm,fill=figorange}]
\foreach \name/\x/\g in {MHA ($H_{\mathrm{kv}}=8$)/0/1, GQA ($H_{\mathrm{kv}}=2$)/40/4, MQA ($H_{\mathrm{kv}}=1$)/80/8} {
  \node[font=\footnotesize\bfseries] at (\x+13.3,22) {\name};
  \foreach \i in {0,...,7} { \node[q] (q\x-\i) at (\x+\i*3.8,12) {}; }
  \pgfmathtruncatemacro{\nk}{8/\g-1}
  \foreach \j in {0,...,\nk} {
    \pgfmathsetmacro{\cx}{\x+(\j*\g+(\g-1)/2)*3.8}
    \node[kv] (k\x-\j) at (\cx,0) {};
    \pgfmathtruncatemacro{\lo}{\j*\g}\pgfmathtruncatemacro{\hi}{\j*\g+\g-1}
    \foreach \i in {\lo,...,\hi} { \draw[black!50] (k\x-\j.north) -- (q\x-\i.south); }
  }
}
\node[lbl,anchor=east] at (-4,12) {Q};
\node[lbl,anchor=east] at (-4,0) {K, V};
\end{tikzpicture}
\caption{クエリのヘッド 8 個（青）と、キー・バリューのヘッド（橙）の対応。線でつながったクエリのヘッドは、同じキーとバリューを使う。キャッシュに保存するのは橙の分だけなので、GQA（中央）ではキャッシュが MHA の 1/4 になる。}
\label{fig:ma-gqa}
\end{figure}

\figref{fig:ma-gqa}で、キャッシュに入るのは下段の橙のマスだけです。
中央の GQA では、4 個のクエリのヘッドが 1 組のキーとバリューを共有するので、キャッシュは MHA の 4 分の 1 になります。
クエリのヘッドは 8 個のままなので、「8 通りの見方で文脈を探す」という能力はある程度保たれます。
Ainslie らは、GQA が MQA に近い速さで MHA に近い品質を出すことを示しました\cite{ainslie2023gqa}。

Llama 3 8B は、クエリのヘッド 32 個に対して $H_{\mathrm{kv}}=8$ を使います。
式\eqref{eq:ma-kv-mem}に入れると $2\times32\times8\times128\times2=131{,}072$ バイト、1 トークンあたり 128 KiB で、Llama 2 7B の 4 分の 1 です。
文脈 8,192 トークンでも 1 GiB で済みます。

\figref{fig:ma-kv-mem}は、同じ式に、いろいろな公開モデルの設定の値を入れたものです（\file{ch11/fig_kv_memory.py}。設定の値は\secref{sec:ma-models}の\tabref{tab:ma-heads}と、本章の本文に挙げたもの）。
本書のミニGPT は 1 トークンあたり 6 KiB（32 ビットで保存した場合）ですが、Llama 3.1 8B は 128 KiB、Llama 3.1 405B は 504 KiB です。
文脈を 131,072 トークンまで使う会話 1 本なら、これに 131,072 を掛けて、Llama 3.1 8B で 16 GiB、405B で 63 GiB になります。
\figref{fig:ma-kvcache}(b) でマスが 1 つ増えるたびに（図は 1 つの層だけですが、実際には全部の層で同時に増えます）、合わせてこれだけのメモリが積み上がっていくのです。
いちばん下の 2 本（Qwen3.5-9B と Qwen3.8-27B）は、線形注意という別の仕組みの層を混ぜたモデルです。
文脈とともにキャッシュが増えるのは、4 層に 1 層の「完全な注意」の層だけなので、層の数のわりに小さくなります（\secref{sec:ma-models}）。

\begin{figure}[tbp]
\centering
\begin{tikzpicture}
\begin{axis}[xbar,width=0.72\linewidth,height=66mm,bar width=2.6mm,table/col sep=comma,
  y dir=reverse,ytick=data,yticklabels from table={code/ch11/out/fig_kv_memory.csv}{name},
  yticklabel style={font=\scriptsize},xmin=0,xmax=600,enlarge y limits=0.06,
  xlabel={1 トークンあたりの KV キャッシュ（KiB）},grid=none,xmajorgrids,
  nodes near coords,nodes near coords align={horizontal},
  nodes near coords style={font=\tiny,text=black}]
\addplot[fill=bookblue!70,draw=none] table[col sep=comma,x=kib,y expr=\coordindex]{code/ch11/out/fig_kv_memory.csv};
\end{axis}
\end{tikzpicture}
\caption{1 トークンあたりの KV キャッシュの大きさ（式\eqref{eq:ma-kv-mem}に設定の値を入れて計算した値。\file{ch11/fig_kv_memory.py}）。本書のミニGPT の 2 つは 32 ビット（4 バイト）、ほかは BF16（2 バイト）で保存した場合。gpt-oss-120b は全 36 層を数えた（半分の層は直近 128 トークンだけを見るので、長い文脈での実際の量はこれより小さい）。Qwen3.5-9B と Qwen3.8-27B は、線形注意の層を除き、完全な注意の層だけを数えた。}
\label{fig:ma-kv-mem}
\end{figure}

\subsection*{実装}
\skimread{押さえる結論は「キーとバリューは $H_{\mathrm{kv}}$ ヘッド分だけ作ってキャッシュに保存し、注意をとる直前に、クエリのヘッドの数まで複製する」です。}

\snippet{GQA（RoPE と KV キャッシュつき）}{code:ma-gqa}{ch11/modern_gpt.py}{gqa}

5 行目の \code{self.dh} は 1 ヘッドの次元 $\dhead$ で、このコードでは第 10 章と同じく $d/H$ です。
実際のモデルでは、$\dhead$ を $d/H$ と別に決めることもあります（\tabref{tab:ma-heads}）。
そのときは $H\dhead\ne d$ になりますが、出力の射影が $d$ 次元に戻すので、ブロックの形は変わりません。
6〜7 行目で、クエリの射影は $H$ ヘッド分、キーとバリューの射影は $H_{\mathrm{kv}}$ ヘッド分の出力を持ちます。
8 行目の \code{self.proj} は、ヘッドの出力をつないだものを $d$ 次元に戻す出力の射影 $\mW_o$ です（第 9 章と同じ）。
キーとバリューは、1 つの線形層 \code{self.kv} でまとめて 1 回の行列積で作ります（2 つの層に分けても結果は同じです）。
12 行目で、クエリを $(B,H,T,\dhead)$ の形にします。
13 行目の \code{view(B, T, 2, self.hkv, self.dh)} でその出力を「キーかバリューか」の軸を含む形に分け、\code{unbind(2)} でその軸に沿って 2 つのテンソルに取り出します。
14 行目で形を $(B,H_{\mathrm{kv}},T,\dhead)$ にそろえます。
15 行目で、クエリとキーにだけ RoPE をかけます。
18〜20 行目が GQA の要で、\code{repeat_interleave} で各キー・バリューのヘッドを \code{rep} $=H/H_{\mathrm{kv}}$ 回ずつ複製し、クエリのヘッドの数にそろえています。
たとえば $H_{\mathrm{kv}}=2$、$H=4$ なら、キーのヘッドは $(\vk^{(0)},\vk^{(0)},\vk^{(1)},\vk^{(1)})$ の順に並び、クエリのヘッド 0 と 1 が $\vk^{(0)}$ を、2 と 3 が $\vk^{(1)}$ を使います（上付きの番号はヘッドの番号）。
\figref{fig:ma-gqa}の、隣どうしのクエリのヘッドが同じ組を使う並びと同じです（\code{repeat} なら $(\vk^{(0)},\vk^{(1)},\vk^{(0)},\vk^{(1)})$ となり、組み方が変わります）。
計算の結果は、グループ内のキーとバリューの重みを同じにしたマルチヘッド注意と同じです。
違いは、パラメータとキャッシュを 1 組分しか持たないことです。
19 行目のコメントの $S$ は、キャッシュの分も含めたキーの数（そこまでの位置の数）です。
複製はキャッシュに保存した\textbf{後}に行うので、保存されるのは $H_{\mathrm{kv}}$ ヘッド分だけです。
21 行目は第 9 章で紹介した \code{F.scaled_dot_product_attention} です。
\code{is_causal=True} は、第 9 章で見たとおり、因果マスクをかけて未来の位置を見せない指定です。
プロンプトを一度に処理するとき（$T=16$ など）は、各位置が自分より後ろを見てはいけないのでマスクが要ります。
生成中に 1 トークンだけを入れるとき（$T=1$）は、過去をすべて見てよいので要りません。
そのため \code{is_causal=(T > 1)} としています。
ただしこの書き方は、「キャッシュがあるときは 1 トークンずつ入れる」使い方を前提にしています。
キャッシュがある状態で 2 トークン以上をまとめて入れると、マスクの位置がずれて正しく計算できません。
22 行目で、ヘッドごとの出力を 1 本のベクトルにつなぎ直し、\code{proj} に通します。

なお、実際の高速な実装は、複製せずに同じキーとバリューを複数のヘッドから読みます。
ここでは式との対応を分かりやすくするために、複製する形で書きました。

\subsection*{FlashAttention：注意の計算そのものを速くする}
\laterread{\secref{sec:gpu-speedup}で高速化の定石を学ぶときに、戻ってくれば足ります。飛ばしても、この章の残りは読めます。}

KV キャッシュと GQA は「何を計算し、何を保存するか」を変える工夫でした。
最後に、「同じ計算をどう実行するか」を変える工夫に一言触れておきます。

第 9 章で見たように、注意は $T\times T$ の重み行列を作るので、長い文脈ではメモリと時間が $T^2$ で増えます。
GPU には、小さくて速いメモリと、大きくて遅いメモリがあります（\chapref{ch:gpu}）。
GPU では、計算そのものよりも、この大きな行列を遅いメモリに書いて読み戻す\textbf{データの移動}に時間がかかります。
FlashAttention\cite{dao2022flashattention}は、クエリとキーを小さなブロックに分け、softmax に要る最大値と合計を、ブロックを読むたびに更新しながら計算することで、$T\times T$ の行列を一度もメモリに書き出さずに、\textbf{数学的に同じ}出力を得ます。
結果は変わらず、速さとメモリだけが改善されるのです。
PyTorch の \code{F.scaled_dot_product_attention} は、GPU では条件が合えばこの種の実装を自動で選びます（\chapref{ch:gpu}）。

この節で分かったのは、KV キャッシュの大きさが式\eqref{eq:ma-kv-mem}で決まり、GQA はキーとバリューのヘッドを共有して、それを $H_{\mathrm{kv}}/H$ 倍に縮めることです。

\section{組み立てる：現代版ミニGPT}\label{sec:ma-assemble}
\skimread{押さえる結論は「部品を入れ替えた現代版ミニGPT は、パラメータが 6\% 少なくても、GPT-2 型と同じくらいには学習できた（検証損失 4.546 と 4.775）。ただし本書の規模では、この差は評価と学習のぶれと同じ程度」です。}

部品がそろったので、組み立てて学習させてみましょう。
この節の問いは、「GPT-2 型と現代版を同じ条件で学習させると、本書の小さな規模で違いは見えるのか」です。

\subsection*{ブロックとモデル}

\snippet{現代版のブロック}{code:ma-block}{ch11/modern_gpt.py}{block}

構造は第 10 章の \code{Block} と同じで、14〜15 行目の 2 本の残差接続も変わりません。
違いは、正規化が RMSNorm に（4、6 行目）、注意が GQA に（5 行目）、FFN が SwiGLU に（11 行目）なったことです。
設定で \code{n_expert} を 1 以上にすると、FFN が次の節で作る MoE に替わります（7〜9 行目）。
モデル全体（\file{ch11/modern_gpt.py} の \code{ModernGPT}）では、位置埋め込みの表がなくなり、代わりに RoPE の角度の表を \code{register_buffer} で持ちます。
\code{register_buffer} は、学習はしないが、モデルと一緒に \code{.to(device)} などで運ばれる値の登録です（第 9 章で因果マスクに使いました）。
表は文脈長の 4 倍の長さまで作ってあるので、学習時より長い系列も計算できます。
\secref{sec:ma-kvcache}の速さの測定で、新しいトークンを 480 個までにしたのも、プロンプトの 16 個と合わせた 496 個を、この表の長さ（$4\times128=512$）に収めるためです。
初期化は第 10 章と同じく標準偏差 0.02 の正規分布で、残差ストリームに足し込む射影（\code{proj} と \code{w_down}）だけを $1/\sqrt{2L}$ 倍に小さくしました。

\subsection*{同じ条件で比べる}

比べる 2 つのモデルの設定は次のとおりです。
どちらも 4 層、4 ヘッド、$d=128$、文脈長 128、語彙 4096（本書の BPE）で、埋め込みと出力の重みは共有します。

\snippet{比較する 2 つのモデル}{code:ma-configs}{ch11/compare_train.py}{configs}

現代版はキー・バリューのヘッドを 2 個（GQA、\code{n_kv_head=2}）、SwiGLU の中間次元を 344 にしました。
パラメータ数は GPT-2 型が 1,329,408、現代版が 1,250,432 です。
現代版のほうが $1{,}329{,}408-1{,}250{,}432=78{,}976$ 少なくなっています。
内訳は 4 つあります。
位置埋め込みの表がなくなります（$128\times128=16{,}384$ の減）。
GQA でキー・バリューの射影が半分になります（1 層あたり $2\times128\times64=16{,}384$、4 層で 65,536 の減）。
正規化のずらし $\bm\beta$ がなくなります（各層に 2 個と最後に 1 個で計 9 個、$9\times128=1{,}152$ の減）。
逆に FFN は少し増えます。
SwiGLU は 1 層 $3\times128\times344=132{,}096$ で、GELU の MLP の $8\times128^2=131{,}072$ より 1,024 多いからです（4 層で 4,096 の増）。
合わせると $16{,}384+65{,}536+1{,}152-4{,}096=78{,}976$ で、差とぴったり合います。
学習は \code{llmbook.train.train_lm} で、800 ステップ、バッチ 16、最大学習率 $2\times10^{-3}$、ウォームアップ 50 ステップです。
データは基盤モデル（\file{base.pt}）と同じ \file{train.bin} と \file{val.bin} を使います。

\begin{figure}[tbp]
\centering
\begin{tikzpicture}
\begin{axis}[xlabel={ステップ},ylabel={検証損失},ymin=4.2,ymax=6,legend pos=north east]
\addplot[bookgray,mark=*,mark size=1.2pt] table[col sep=comma,x=step,y=val_loss]{code/ch11/out/compare_gpt2.csv};
\addlegendentry{GPT-2 型（1,329,408）}
\addplot[bookaccent,mark=square*,mark size=1.2pt] table[col sep=comma,x=step,y=val_loss]{code/ch11/out/compare_modern.csv};
\addlegendentry{現代版（1,250,432）}
\end{axis}
\end{tikzpicture}
\caption{GPT-2 型と現代版ミニGPT の検証損失（800 ステップ、各 1 回の学習）。ステップ 0 の 8.3 は縦軸の範囲外。検証損失は、毎回ランダムに選んだ 20 バッチで見積もっているので、曲線は上下に揺れている。}
\label{fig:ma-compare}
\end{figure}

\figref{fig:ma-compare}が結果です。
最後（ステップ 800）の検証損失は、GPT-2 型が 4.775、現代版が 4.546 でした。
ステップ 200 以降の 7 点のうち 5 点で、現代版のほうが低くなっています。
ただし、この比較から「現代版が 0.2 良い」と結論するのは早すぎます。
曲線は 100 ステップごとに 0.2〜0.3 も上下しているからです。
この揺れには、2 つの原因が重なっています。
1 つは評価のぶれで、検証損失を毎回ランダムに選んだ 20 バッチで見積もっていることによります（評価に使うバッチを固定するか増やせば減ります）。
もう 1 つは学習のぶれで、初期値やバッチの順番の乱数を変えると、学習の結果そのものが変わることによります。
学習は 1 回ずつなので、0.2 の差がどちらによるものかは、この実験からは区別できません。
言えるのは、「パラメータが 6\% 少ない現代版が、少なくとも同じくらいには学習できた」ことまでです。
大規模なモデルで現代の部品が優れていることは、多くの比較実験で確かめられています\cite{shazeer2020glu,su2024roformer,touvron2023llama}。
しかし本書の規模の CPU の数分の学習で、その差をはっきり示すのは困難です。

現代版の利点がはっきり出るのは、むしろ損失以外の面です。
現代版は位置の表を持たないので、学習時より長い文脈も計算できます。
KV キャッシュのメモリは GQA で半分です（1 トークンあたり 4 KiB が 2 KiB に。32 ビットで保存した場合）。
これらは小さなモデルでも確かめられる、設計そのものの性質です。

この節で分かったのは、現代版ミニGPT は部品を入れ替えるだけで組み立てられ、本書の規模では損失の差は、評価と学習のぶれと同じ程度だということです。

\section{MoE：1 トークンの計算を増やさずにパラメータを増やす}\label{sec:ma-moe}

最後の差分は、FFN の「大きさ」の扱いです。
2025 年以降の大きな公開モデルの多くは、\term{混合エキスパート}{mixture of experts; MoE}という形の FFN を使っています。
この節の問いは、「パラメータ数と 1 トークンあたりの計算量を切り離すには、どうすればよいか」です。

\subsection*{何に困っているか}

スケーリング則（\chapref{ch:pretraining}で詳しく扱います）によれば、十分なデータがあれば、モデルのパラメータを増やすほど損失は下がります。今はこの一点だけを使います。
これまでの、すべてのトークンがすべてのパラメータを通るモデルを、MoE と区別して\textbf{密な}（dense）モデルと呼びます。
ところが密なモデルでは、パラメータを 2 倍にすると、1 トークンを処理する計算量もほぼ 2 倍になります。
すべてのトークンが、すべてのパラメータを通るからです。
学習の費用も、生成の費用も、パラメータ数に比例して増えてしまいます。

ここで、人間の組織を考えてみましょう。
大きな病院には、多くの専門医がいます。
しかし一人の患者を診るのは、そのうちの数人だけです。
病院全体の知識は医師の人数に比例して増えますが、一人の患者にかかる手間は増えません。
受付が症状を見て、適切な科に振り分けるからです。

MoE は、FFN をこの病院のように作ります。
FFN を $E$ 個の\textbf{エキスパート}（expert、それぞれが普通の FFN）に分け、\textbf{ルータ}（router、受付）が各トークンを $k$ 個のエキスパートにだけ送ります。
たとえの対応は、専門医がエキスパート、受付がルータ、患者がトークンです。
この $k$ は、\chapref{ch:lm}の top-$k$ サンプリングと同じく「上位 $k$ 個」の意味ですが、選ぶ対象はトークンではなくエキスパートです。
また、エキスパートの数 $E$ は細字で、トークン埋め込みの行列 $\mE$（太字）とは別のものです。
パラメータは $E$ 個分ありますが、1 トークンが通るのは $k$ 個分だけです。
$E$ をいくら増やしても、1 トークンの FFN の計算は $k$ 個分のままです。
同じ大きさの FFN を 1 個持つ密なモデルと比べると、top-1 なら同じ、top-2 なら 2 倍です。
節の題の「計算を増やさずに」は、「$E$ を増やしても、1 トークンの計算が増えない」という意味です。
ただし、このたとえには限界があります。
実際のエキスパートは「数学の専門家」「文法の専門家」のように人間に分かる分担をしているとは限らず、どう分担するかは学習の結果として決まります。

\subsection*{ルータと top-$k$ ゲーティング}

1 つの位置の入力 $\vx\in\R^d$ に対し、ルータは線形変換 $\mW_r\in\R^{d\times E}$（ルータの重み）と softmax で、各エキスパートへの「適性」を確率として出します。
この確率のベクトル $\vp$ を、\textbf{ルータの確率}と呼びます（位置埋め込みの $\vp_t$ とは別のものです）。
\begin{equation}
\vp=\softmax(\mW_r^\T\vx)\in\R^E .
\label{eq:ma-router}
\end{equation}
このうち確率の大きい上位 $k$ 個のエキスパートの番号の集合を $\mathrm{Top}_k$ とし、その確率を足して 1 になるように割り直したものを重み $g_e=p_e/\sum_{j\in\mathrm{Top}_k}p_j$ とします。
これは、上位 $k$ 個のロジットだけで softmax をとり直した値と同じです（softmax の分母が打ち消し合うからです）。
$k=1$ のときは、選ばれるのは 1 個なので、重みはいつも $g_e=p_e/p_e=1$ になります（この点は、後の実験で大事になります）。
MoE の出力は、選ばれたエキスパートの出力の重み付き和です。
\begin{equation}
\mathrm{MoE}(\vx)=\sum_{e\in\mathrm{Top}_k} g_e\,\mathrm{FFN}_e(\vx).
\label{eq:ma-moe}
\end{equation}
選ばれなかったエキスパートは計算しません。
この「上位 $k$ 個だけを使う」仕組みを、\term{top-$k$ ゲーティング}{top-$k$ gating}と呼びます\cite{shazeer2017outrageously}。
ここでのゲーティングは、エキスパートを選んで重みを付けることです。
SwiGLU のゲート（成分ごとの倍率）とは別の使い方です。

\begin{example}[4 エキスパート・top-2]
$E=4$ で、あるトークンのルータの出力が $\vp=(0.1,\,0.5,\,0.3,\,0.1)$ だったとします。
エキスパートの番号は、コードと同じく 0 から数えます。
上位 2 個はエキスパート 1 と 2 なので、$g_1=0.5/0.8=0.625$、$g_2=0.3/0.8=0.375$ です。
出力は $0.625\,\mathrm{FFN}_1(\vx)+0.375\,\mathrm{FFN}_2(\vx)$ で、エキスパート 0 と 3 は計算されません。
\end{example}

\begin{figure}[tbp]
\centering
\begin{tikzpicture}[x=1mm,y=1mm,e/.style={box,font=\footnotesize,minimum width=15mm,minimum height=9mm,align=center}]
\node[box,font=\footnotesize] (x) at (0,0) {$\vx$（1 トークン）};
\node[purplebox,font=\footnotesize,align=center] (r) at (-52,0) {ルータ\\$\vp=\softmax(\mW_r^\T\vx)$};
\draw[arr] (x.west) -- (r.east);
\node[e,fill=figgray,draw=black!40,text=black!50] (e1) at (-36,24) {FFN$_0$\\[-0.5mm]{\scriptsize\color{bookpurple}$p_0=0.1$}};
\node[e] (e2) at (-12,24) {FFN$_1$\\[-0.5mm]{\scriptsize\color{bookpurple}$p_1=0.5$}};
\node[e] (e3) at (12,24) {FFN$_2$\\[-0.5mm]{\scriptsize\color{bookpurple}$p_2=0.3$}};
\node[e,fill=figgray,draw=black!40,text=black!50] (e4) at (36,24) {FFN$_3$\\[-0.5mm]{\scriptsize\color{bookpurple}$p_3=0.1$}};
\node[op] (s) at (0,42) {$\Sigma$};
\draw[darr,black!30] (x.north) -- (e1.south);
\draw[arr] (x.north) -- (e2.south);
\draw[arr] (x.north) -- (e3.south);
\draw[darr,black!30] (x.north) -- (e4.south);
\draw[bookpurple,densely dashed,thick] (r.north) |- node[pos=0.3,font=\scriptsize,text=bookpurple,left,align=right] {確率と\\選択} (32,32);
\foreach \ee in {e1,e2,e3,e4} {\draw[->,bookpurple,densely dashed,thick] ([xshift=-4mm]\ee.north |- 0,32) -- ([xshift=-4mm]\ee.north);}
\draw[arr,preaction={draw=white,line width=2.4pt,-}] (e2.north) -- node[lbl,above left,pos=0.72,inner sep=1pt] {$\times0.625$} (s);
\draw[arr,preaction={draw=white,line width=2.4pt,-}] (e3.north) -- node[lbl,above right,pos=0.72,inner sep=1pt] {$\times0.375$} (s);
\draw[arr] (s) -- ++(0,9);
\end{tikzpicture}
\caption{MoE の層（$E=4$、top-2）。黒の矢印は、1 トークンの入力 $\vx$ がエキスパートに入る流れで、各エキスパートの入力は同じ $\vx$。紫はルータで、各エキスパートの確率（箱の中の紫の数）を出し、上位 2 個（エキスパート 1 と 2）を選ぶ。紫の破線は、その確率と選んだ結果を 4 個のエキスパートすべてに渡すことを表す（ルータの出力がエキスパートの入力になるのではない）。選ばれた 2 個だけを計算し、割り直した重みを掛けて足し合わせる。灰色のエキスパート（破線の矢印）は、このトークンでは計算されない。}
\label{fig:ma-moe}
\end{figure}

\figref{fig:ma-moe}は上の数値例を図にしたものです。
トークンごとに選ばれるエキスパートは変わるので、同じ文の中でも、あるトークンは 1 と 2、次のトークンは 0 と 3、ということが起きます。

\subsection*{負荷分散：一部のエキスパートに偏らせない}

MoE をそのまま学習させると、困ったことが起きやすくなります。
ルータが学習で更新されるとき、学習の初めに少し多く選ばれたエキスパートは、多くのトークンで学習するので、そのトークンを少しうまく処理できるようになります。
すると、ルータはそのエキスパートをさらに多く選ぶようになります。
この自己強化の循環が続くと、最後には一部のエキスパートだけが使われ、残りはほとんど使われなくなります。
これではパラメータを増やした意味がありません。
しかも複数の GPU にエキスパートを分けて置く実際の学習では、トークンが集中したエキスパートの GPU の処理を、ほかの GPU が待つことになります。

そこで、偏りが大きいほど大きくなる\term{負荷分散損失}{load-balancing loss}を、言語モデルの損失に小さな係数で足します。
本来の損失に補助として足す損失（auxiliary loss）なので、\textbf{補助損失}とも呼びます。
本章では、2 つの呼び名を同じ意味で使います。
これまでの次のトークン予測の損失を $\calL_{\mathrm{LM}}$、負荷分散損失を $\calL_{\mathrm{aux}}$（aux は auxiliary の略）と書くと、学習で下げるのは $\calL_{\mathrm{LM}}+c\,\calL_{\mathrm{aux}}$ です。
係数 $c$ は 0.01 のような小さな値にして、本来の目的である $\calL_{\mathrm{LM}}$ を妨げすぎないようにします。

議論の骨組みを先に示します。
(1) 偏りの大きさを、エキスパートごとの「1 位に選ばれた割合」と「確率の平均」の積の和で測ります。
(2) 選ばれた割合は数えるだけの量なので、定数とみなします。
(3) すると損失は確率の平均の 1 次式になり、混んだエキスパートの確率ほど強く押し下げる勾配が出ます。
Switch Transformer\cite{fedus2022switch}の形は次のとおりです。
バッチ内の $N=B\times T$ 個のトークンに $1,\dots,N$ の通し番号 $t$ を付けて、
\begin{equation}
\calL_{\mathrm{aux}}=E\sum_{e=0}^{E-1}f_e\,P_e,\qquad
f_e=\frac{1}{N}\sum_{t=1}^{N}\mathbf{1}[\text{トークン $t$ の 1 位が $e$}],\quad
P_e=\frac{1}{N}\sum_{t=1}^{N}p_{t,e}.
\label{eq:ma-aux}
\end{equation}
$f_e$ はエキスパート $e$ が実際に 1 位に選ばれた割合、$P_e$ はルータが $e$ に与えた確率の平均です。
$p_{t,e}$ は、トークン $t$ のルータの確率 $\vp$ の、エキスパート $e$ の成分です。
$\mathbf{1}[\cdot]$ は、中の条件が成り立てば 1、そうでなければ 0 を表します。
top-2 以上で使うときも、この形では 1 位だけを数えます（本章のコードも同じです）。

\begin{figure}[tbp]
\centering
\begin{tikzpicture}[x=1mm,y=1mm,
  pc/.style={cell,minimum width=13mm,minimum height=5.2mm,font=\scriptsize},
  tc/.style={pc,fill=figorange,font=\scriptsize\bfseries}]
\foreach \e in {0,...,3} { \node[font=\scriptsize] at (\e*14,7) {$e=\e$}; }
\node[font=\scriptsize,anchor=east] at (-8,7) {エキスパート};
\foreach \t/\pa/\pb/\pc/\pd/\top in {1/0.1/0.5/0.3/0.1/1, 2/0.6/0.2/0.1/0.1/0, 3/0.2/0.4/0.3/0.1/1, 4/0.1/0.2/0.2/0.5/3} {
  \node[font=\scriptsize,anchor=east] at (-8,-\t*6+6) {トークン \t\ の $\vp$};
  \foreach \e/\v in {0/\pa,1/\pb,2/\pc,3/\pd} {
    \ifnum\e=\top \node[tc] at (\e*14,-\t*6+6) {\v};
    \else \node[pc] at (\e*14,-\t*6+6) {\v}; \fi
  }
}
\draw[black!50] (-36,-21.5) -- (49,-21.5);
\foreach \lab/\a/\b/\c/\d/\y in {
  {1 位の回数}/1/2/0/1/-25,
  {$f_e$（回数 $\div N$）}/0.25/0.5/0/0.25/-31,
  {$P_e$（列の平均）}/0.25/0.325/0.225/0.2/-37,
  {$f_e P_e$}/0.0625/0.1625/0/0.05/-43,
  {勾配 $E f_e$}/1/2/0/1/-49} {
  \node[font=\scriptsize,anchor=east] at (-8,\y) {\lab};
  \foreach \e/\v in {0/\a,1/\b,2/\c,3/\d} { \node[font=\scriptsize] at (\e*14,\y) {\v}; }
}
\node[font=\scriptsize] at (21,-56) {$f_eP_e$ の和 0.275 $\to$ $\calL_{\mathrm{aux}}=4\times0.275=1.10$};
\end{tikzpicture}
\caption{負荷分散損失（式\eqref{eq:ma-aux}）を、$N=4$ トークン・$E=4$ エキスパートの小さな例で計算する（数値は説明のための例）。上の 4 行は各トークンのルータの確率 $\vp$ で、橙のマスが 1 位。トークン 1 は\figref{fig:ma-moe}と同じ確率。列ごとに橙を数えて $N$ で割ったものが $f_e$、列の平均が $P_e$。最後の行が $P_e$ についての勾配で、混んだエキスパート 1 の確率が最も強く押し下げられる。}
\label{fig:ma-aux}
\end{figure}

\figref{fig:ma-aux}は、4 トークン・4 エキスパートの小さな例で、$f_e$ と $P_e$ を作るところを示したものです。
各行が 1 トークンのルータの確率で、橙のマスが 1 位です。
1 位はトークンごとに違い、この例ではエキスパート 1 が 2 回、0 と 3 が 1 回ずつ、2 は 0 回です。
列ごとに橙のマスを数えて $N=4$ で割ったものが $f_e$、列の平均が $P_e$ で、$\calL_{\mathrm{aux}}=4\times0.275=1.10$ になります。
両方が均等（どちらも $1/E$）なら、$\calL_{\mathrm{aux}}=E\cdot E\cdot(1/E)^2=1$ です。
1 つのエキスパートに集中すると（$f_e=P_e=1$）、$\calL_{\mathrm{aux}}=E$ と大きくなります。
たとえば $E=2$ で $f=(0.9,0.1)$、$P=(0.8,0.2)$ と偏っていれば、$2\times(0.9\times0.8+0.1\times0.2)=1.48$ です。
ただし 1 は目安で、最小値とは限りません。
$f_e$ と $P_e$ は別の量なので、1 を下回ることもあります。
たとえば 3 トークン・2 エキスパートで、$\vp$ が $(0.51,0.49)$、$(0.51,0.49)$、$(0,1)$ なら、$f=(2/3,1/3)$、$P=(0.34,0.66)$ で、$\calL_{\mathrm{aux}}\approx0.893$ です。
大事なのは値の大小ではなく、次に見る勾配の向きです。

$f_e$ は数を数えるだけなので微分できず、定数として扱います。
一方 $P_e$ はルータの出力の平均なので微分でき、
\[
\pd{\calL_{\mathrm{aux}}}{P_e}=E f_e
\]
です。
混んでいるエキスパート（$f_e$ が大きい）ほどこの勾配が大きいので、勾配降下法で損失を下げると、そのエキスパートの確率 $P_e$ が、ほかのエキスパートと比べて強く押し下げられます。
確率の合計はいつも 1 なので、下がった分はほかのエキスパートに移ります。
\figref{fig:ma-aux}の例なら勾配は $(1,2,0,1)$ で、2 回選ばれたエキスパート 1 の確率が最も強く下げられます。
1 回も選ばれなかったエキスパート 2 の確率は、この項からは下げられず、ほかが下がる分だけ相対的に上がります。
これが「混んだエキスパートを避けさせる」仕組みです。

\subsection*{実装}
\skimread{押さえる結論は「エキスパートごとに、自分を選んだトークンだけを計算し、重みを掛けて足す。式\eqref{eq:ma-router}〜\eqref{eq:ma-aux}が、そのまま 20 行ほどのコードになる」です。}

\snippet{小さな MoE 層}{code:ma-moe}{ch11/moe.py}{moe}

5 行目がルータ $\mW_r$、6〜7 行目が $E$ 個のエキスパートで、それぞれが\secref{sec:ma-swiglu}の SwiGLU です。
13 行目で $(B,T,C)$ を $(N,C)$ の行列にし、$N=B\times T$ 個のトークンを、1 トークン 1 行で縦に並べます。
14 行目が式\eqref{eq:ma-router}、15〜16 行目が上位 $k$ 個の選択と重みの割り直しです。
\code{topk} は値 \code{w} と番号 \code{idx} を返し、どちらも形は $(N,k)$ です。
18〜21 行目がエキスパートごとのループです。
\code{(idx == e).nonzero(as_tuple=True)} は、エキスパート \code{e} を選んだ行の番号 \code{tok} と、それが何番目の候補だったか \code{slot} を返します。
そのトークンだけをエキスパートに通し、重みを掛けて \code{y}（17 行目で 0 に初期化した、\code{x} と同じ形の入れ物）に足し込みます。
\code{w[tok, slot, None]} の \code{None} は長さ 1 の軸を足して形を $(n,1)$ にする書き方で、ブロードキャストで $(n,C)$ の出力の各行に重みを掛けるためのものです。
23〜25 行目が式\eqref{eq:ma-aux}です。
\code{F.one_hot(idx[:, 0], ...)} が 1 位のエキスパートを 1 にしたベクトルで、その平均が $f_e$、\code{probs.mean(0)} が $P_e$ です。
26 行目では、後で偏りを調べるために、各エキスパートが選ばれた回数を記録しています。
損失と回数は層の属性（8〜9 行目で用意した \code{aux_loss} と \code{usage}）に残すので、学習ループから後で読み出せます。
27 行目で、形を元の $(B,T,C)$ に戻します。

実際の大規模な実装では、このループの代わりに、トークンをエキスパートごとに並べ替えてまとめて行列積をする工夫や、各エキスパートが受け付けるトークン数に上限（capacity）を設ける工夫が入ります。

\subsection*{実験：負荷分散損失の効果}
\skimread{押さえる結論は「負荷分散損失がないと、4\% しか使われないエキスパートが現れる。足すと、どれも 13\% 以上使われる。ただし本書の規模では、MoE の検証損失（4.896。補助損失なしで 4.800）は、密なモデル（4.793）より良くならなかった」です。}

小さなモデル（2 層、$d=96$）で、次の 3 つを同じ条件（800 ステップ、バッチ 16、系列長 64）で学習させます。
\begin{itemize}
\item \textbf{密}: 普通の SwiGLU（中間次元 256）。
\item \textbf{MoE}: 同じ大きさの SwiGLU を 4 個のエキスパートにし、top-1 で選ぶ。負荷分散損失の係数 0.01。
\item \textbf{MoE（補助損失なし）}: 上と同じで、負荷分散損失を足さない。
\end{itemize}
top-1 なので、1 トークンあたりの FFN の計算量は 3 つとも同じです（ルータの分だけ MoE がわずかに多い）。

ここで、1 つ注意があります。
top-1 では、選ばれたエキスパートの重みはいつも $g_e=1$ でした（式\eqref{eq:ma-moe}の前）。
MoE の出力はルータの確率の値によらないので、言語モデルの損失 $\calL_{\mathrm{LM}}$ からは、ルータの重み $\mW_r$ に勾配が届きません。
手元で確かめると、$\calL_{\mathrm{LM}}$ からのルータの勾配の最大値は約 $6\times10^{-11}$ で、丸め誤差の大きさでした（同じ初期値で top-2 にすると約 $1\times10^{-3}$ になります）。
したがって、この実験のルータは、補助損失ありでは負荷分散損失 $\calL_{\mathrm{aux}}$ だけから学習し、補助損失なしでは、ほぼ初期値のまま変わりません。
同じ学習を手元で再現して確かめると、補助損失なしのルータの重みは、初期値とのコサイン類似度（\chapref{ch:math}）が 0.9999 で、重み減衰（\chapref{ch:training}）で大きさが約 0.92 倍に縮んだだけでした（補助損失ありでは 0.80〜0.87）。
Switch Transformer の top-1 は、割り直す前の確率 $p_e$ をそのまま出力に掛けるので、ルータにも $\calL_{\mathrm{LM}}$ の勾配が届きます\cite{fedus2022switch}。
本章のコードは、top-2 以上のための割り直しを top-1 にもそのまま使ったので、この違いが生じました。
ルータが $\calL_{\mathrm{LM}}$ からも学習する形は、演習 6 の top-2 で試せます。
一方、FFN のパラメータは MoE が 4 倍です。
SwiGLU 1 個は $3\times96\times256=73{,}728$ 個で、MoE では各層に 3 個分が増えます（2 層で $442{,}368$ 個）。
ルータ（$96\times4=384$ 個）も 2 層で 768 個増えるので、総パラメータは $596{,}448+442{,}368+768=1{,}039{,}584$ になります。

\snippet{MoE の実験の学習ループ（抜粋）}{code:ma-moe-loop}{ch11/moe_experiment.py}{moe-loop}

学習ループは第 10 章の \code{train_lm} とほぼ同じで、違いは 16〜17 行目で負荷分散損失を足しているところです。
4〜6 行目では、直前の学習ステップで最も多くのトークンを受け持ったエキスパートの割合を記録しています。
均等なら 0.25 です。

\begin{table}[tbp]
\centering
\caption{MoE の実験の結果（\file{ch11/out/moe_*.txt}）。エキスパートの割合は、学習後に検証データの 1 バッチ（1,024 トークン）で各エキスパートが 1 位に選ばれたトークンの割合。均等なら 0.25。割合はエキスパート 0〜3 の順。}
\label{tab:ma-moe}
\small
\begin{tabular}{@{}lrrll@{}}
\toprule
モデル & 総パラメータ & 検証損失 & 層 0 の割合 & 層 1 の割合 \\
\midrule
密（SwiGLU） & 596,448 & 4.793 & --- & --- \\
MoE（補助損失 0.01） & 1,039,584 & 4.896 & 0.30 0.23 0.13 0.34 & 0.31 0.24 0.22 0.22 \\
MoE（補助損失なし） & 1,039,584 & 4.800 & 0.11 0.34 0.39 0.16 & 0.04 0.41 0.40 0.15 \\
\bottomrule
\end{tabular}
\end{table}

\tabref{tab:ma-moe}が結果です。
検証損失は、学習の最後に、検証データの 100 バッチを使って測った値です（途中の記録の 16 バッチより多く使い、ぶれを小さくしました）。
途中の記録をとる 7 行目の \code{estimate_loss(model, val, 16)} の 16 はバッチの大きさで、バッチの数は既定の 20 です。
エキスパートの割合は、\figref{fig:ma-moe-share}に棒グラフでも示しました。

\begin{figure}[tbp]
\centering
\begin{tikzpicture}
\begin{axis}[width=0.5\linewidth,height=40mm,ybar,bar width=3.2mm,ymin=0,ymax=0.62,ytick={0,0.2,0.4,0.6},xmin=-0.6,xmax=3.6,
  xtick={0,1,2,3},xlabel={エキスパートの番号},ylabel={選ばれた割合},grid=none,ymajorgrids,
  legend pos=north west,legend style={font=\scriptsize},legend columns=2,
  title={\footnotesize (a) 補助損失なし},title style={yshift=-1mm}]
\addplot[fill=bookblue,draw=none] coordinates {(0,0.11) (1,0.34) (2,0.39) (3,0.16)};
\addlegendentry{層 0}
\addplot[fill=bookaccent,draw=none] coordinates {(0,0.04) (1,0.41) (2,0.40) (3,0.15)};
\addlegendentry{層 1}
\draw[black!70,dashed,thick] (axis cs:-0.6,0.25) -- (axis cs:3.6,0.25);
\end{axis}
\end{tikzpicture}\hfill
\begin{tikzpicture}
\begin{axis}[width=0.5\linewidth,height=40mm,ybar,bar width=3.2mm,ymin=0,ymax=0.62,ytick={0,0.2,0.4,0.6},xmin=-0.6,xmax=3.6,
  xtick={0,1,2,3},xlabel={エキスパートの番号},ylabel={選ばれた割合},grid=none,ymajorgrids,
  legend pos=north west,legend style={font=\scriptsize},legend columns=2,
  title={\footnotesize (b) 補助損失あり（係数 0.01）},title style={yshift=-1mm}]
\addplot[fill=bookblue,draw=none] coordinates {(0,0.30) (1,0.23) (2,0.13) (3,0.34)};
\addlegendentry{層 0}
\addplot[fill=bookaccent,draw=none] coordinates {(0,0.31) (1,0.24) (2,0.22) (3,0.22)};
\addlegendentry{層 1}
\draw[black!70,dashed,thick] (axis cs:-0.6,0.25) -- (axis cs:3.6,0.25);
\end{axis}
\end{tikzpicture}
\caption{各エキスパートが 1 位に選ばれたトークンの割合（\tabref{tab:ma-moe}の数値）。破線は均等な場合の 0.25。(a) 補助損失なしでは、エキスパート 1 と 2 に集中し、層 1 のエキスパート 0 はほとんど使われない。(b) 補助損失を足すと、偏りが小さくなる（層 1 は 0.22〜0.31。層 0 はまだ 0.13〜0.34 の幅がある）。青が層 0、赤が層 1（(a)(b) 共通）。}
\label{fig:ma-moe-share}
\end{figure}

まず、エキスパートの使われ方を見ましょう。
補助損失なしでは、層 1 のエキスパート 0 が 4\% しか使われず、エキスパート 1 と 2 が 4 割ずつを受け持っています。
ただし上で注意したとおり、補助損失なしのルータはほぼ初期値のままなので、この偏りは、\textbf{ルータが学習して}生じたものではありません。
学習で各位置のベクトル（ルータへの入力）が変わっていき、初期値のままのルータがそれを偏って振り分けた結果だと考えられます。
何も手を打たなければ、800 ステップの短い学習でもこれだけ偏ることがある、と読むのが正確です。
補助損失ありでは、最も少ないエキスパートでも 13\%、層 1 では 22〜31\% と、均等の 25\% に近づきました。
式\eqref{eq:ma-aux}の補助損失が、狙いどおりに働いています。

次に損失です。
ここは正直に書く必要があります。
MoE は FFN のパラメータが 4 倍あるのに、検証損失は密なモデル（4.793）より良くなりませんでした（4.896 と 4.800）。
理由として考えられるのは、各エキスパートが全トークンの約 4 分の 1 しか見ないので、800 ステップでは一つひとつのエキスパートの学習が足りないことです。
補助損失ありのほうが損失が少し高い理由は、この実験からは言えません。
ルータは $\calL_{\mathrm{LM}}$ から学習していないので、「言語モデルにとって良い振り分け」を学ぶ仕組みが、もともとないからです。
いずれにせよ、この規模の 1 回ずつの比較では、0.1 程度の差は誤差の範囲だと考えるべきです。

大規模では事情が違います。
Switch Transformer は、同じ計算量の密なモデルと比べて、同じ損失に達するまでの学習を数倍速くできたと報告しています\cite{fedus2022switch}。
Mixtral 8x7B は、総パラメータ 467 億のうち 1 トークンあたり 129 億しか使わずに、当時の 700 億パラメータ級の密なモデルに匹敵する性能を示しました\cite{jiang2024mixtral}。
MoE の利点は、十分なデータと学習量があって初めて引き出せるものなのです。

\subsection*{総パラメータとアクティブパラメータ}

MoE のモデルを語るときは、2 つのパラメータ数を区別します。
\textbf{総パラメータ}は、すべてのエキスパートを含めた数で、メモリに置く必要がある量です。
\textbf{アクティブパラメータ}は、1 トークンの計算で実際に使われる数で、計算量（したがって速さと費用）を決める量です。
たとえば Qwen3-30B-A3B という名前は、「総パラメータ 300 億、アクティブ 30 億」を表しています（A は active）。
128 個のエキスパートから 8 個を選ぶので、FFN のうち使われるのは 16 分の 1 だけです\cite{yang2025qwen3}。

DeepSeek-V3\cite{liu2024deepseekv3}は、総パラメータ 6,710 億、アクティブ 370 億です。
256 個の「ルーティングされる」エキスパートから 8 個を選び、さらに全トークンが必ず通る\textbf{共有エキスパート}を 1 個持ちます。
負荷分散は、補助損失に頼る代わりに、各エキスパートのルータのスコアに足すずらしの値を、混み具合に応じて少しずつ調整する方法を主に使います（補助損失なしの負荷分散）。
補助損失が言語モデルの学習を妨げる、という問題を避けるための工夫です。
本章の実験の 4.896 と 4.800 の差（\tabref{tab:ma-moe}）は誤差の範囲で、しかもルータが $\calL_{\mathrm{LM}}$ から学習しない設定なので、この問題の裏付けにはなりません。

負荷分散のやり方は、モデルによって選択が分かれています。
\tabref{tab:ma-balance}に、本章と\tabref{tab:ma-models}のモデルに出てくる方法を並べました。
補助損失を使うモデル（Mixtral、Qwen3）と使わないモデル（DeepSeek-V3）が並んでいて、どれが良いかは決着していません。
補助損失は簡単でどんなルータにも使えますが、$\calL_{\mathrm{LM}}$ と別の目的を足すので、係数が大きいと言語モデルの学習を妨げるおそれがあります。
補助損失を使わない方法は、その心配を避けるかわりに、ずらしの値を調整する規則を別に決める必要があります。

\begin{table}[tbp]
\centering
\caption{MoE の負荷分散の方法の比較（事実は本章の本文と、各モデルの公開された設定・技術報告による）。}
\label{tab:ma-balance}
\footnotesize
\setlength{\tabcolsep}{3pt}
\begin{tabular}{@{}>{\raggedright\arraybackslash}p{21mm}>{\raggedright\arraybackslash}p{30mm}>{\raggedright\arraybackslash}p{21mm}>{\raggedright\arraybackslash}p{24mm}>{\raggedright\arraybackslash}p{24mm}@{}}
\toprule
方法 & 仕組みの要点 & 長所 & 短所・費用 & 使っている例 \\
\midrule
補助損失（負荷分散損失） & 偏りが大きいほど大きくなる $\calL_{\mathrm{aux}}$ を小さな係数で $\calL_{\mathrm{LM}}$ に足す（式\eqref{eq:ma-aux}） & 簡単。どのルータにも使える & 係数が大きいと $\calL_{\mathrm{LM}}$ の学習を妨げる & Switch Transformer、Mixtral（係数 0.02）、Qwen3（バッチ全体で数える形。235B-A22B は係数 0.001）、本章の実験 \\
補助損失なし（ずらしの調整） & 各エキスパートのスコアに足すずらしの値を、混み具合に応じて少しずつ調整する & 学習の目的（損失）を変えない & 調整の規則と速さを別に決める必要がある & DeepSeek-V3 \\
ハッシュによる振り分け & 学習せず、決まった規則（トークンの番号から計算する値）で行き先を決める & 偏りが学習で生じない & 行き先が文脈によらない & DeepSeek-V4（最初の 3 つの MoE 層） \\
\bottomrule
\end{tabular}
\end{table}


この節で分かったのは、MoE がルータで各トークンを一部のエキスパートにだけ送ることで、1 トークンの計算量を増やさずにパラメータを増やす仕組みであり、偏りを防ぐ負荷分散が欠かせないこと、その方法にはまだ複数の選択肢があることです。

\section{実在のモデルの設定ファイルを読む}\label{sec:ma-models}
\skimread{押さえる結論は「設定ファイルの数個の値（層の数、次元、ヘッドの数、FFN の中間次元）から、パラメータ数と KV キャッシュの大きさを自分で計算できる」です。表とコードには、\chapref{ch:infra}で戻ってきます。}

ここまでの部品で、公開されている LLM の設定ファイルがほとんど読めるようになりました。
この節の問いは、「Llama 3 や DeepSeek-V3 の設定を見て、何が分かるのか。パラメータ数やメモリは自分で計算できるのか」です。

\subsection*{代表的なモデルの構成}

\tabref{tab:ma-models}に、代表的な公開モデルの構成をまとめました。
数値は各モデルの公開された設定ファイルと技術報告によります\cite{radford2019language,dubey2024llama3,jiang2024mixtral,yang2025qwen3,liu2024deepseekv3}。
Gemma 4 と gpt-oss は、モデルの公開ページに置かれた設定ファイルの値です。
表の Llama 3.1 8B は、文脈長（Llama 3 は 8,192）を除いて、Llama 3 8B と同じ構成です。
本章のほかの場所にある「Llama 3 8B」の計算は、Llama 3.1 8B にもそのまま当てはまります。

\begin{table}[tbp]
\centering
\caption{代表的な公開モデルの構成。「Q/KV」はクエリとキー・バリューのヘッド数（等しければ MHA、KV が少なければ GQA）。MoE の行の FFN は「エキスパート数・選ぶ数（1 個の中間次元）」で、DeepSeek-V3 の「＋1」は共有エキスパート。MLA は表の後で説明する。文脈長は扱える最大の長さで、K は 1,024 を表す（128K は $2^{17}=131{,}072$、約 13 万）。}
\label{tab:ma-models}
\footnotesize
\setlength{\tabcolsep}{3.6pt}
\begin{tabular}{@{}lrrrllr@{}}
\toprule
モデル（公開年） & 総／アクティブ & $L$ & $d$ & Q/KV & FFN & 文脈長 \\
\midrule
GPT-2（2019） & 1.24 億 & 12 & 768 & 12/12 & 3072, GELU & 1,024 \\
Llama 3.1 8B（2024） & 80 億 & 32 & 4096 & 32/8 & 14336, SwiGLU & 131,072 \\
Llama 3.1 405B（2024） & 約 4,050 億 & 126 & 16384 & 128/8 & 53248, SwiGLU & 131,072 \\
Mixtral 8x7B（2023） & 467 億／129 億 & 32 & 4096 & 32/8 & 8・2（14336） & 32,768 \\
Qwen3-8B（2025） & 82 億 & 36 & 4096 & 32/8 & 12288, SwiGLU & 約 128K \\
Qwen3-235B-A22B（2025） & 2,350 億／220 億 & 94 & 4096 & 64/4 & 128・8（1536） & 約 128K \\
DeepSeek-V3（2024） & 6,710 億／370 億 & 61 & 7168 & MLA & 256・8＋1（2048） & 約 128K \\
gpt-oss-120b（2025） & 1,168 億／51 億 & 36 & 2880 & 64/8 & 128・4（2880） & 131,072 \\
Gemma 4 12B（2026） & 120 億 & 48 & 3840 & 16/8 & 15360, GeGLU & 262,144 \\
\bottomrule
\end{tabular}
\end{table}

\tabref{tab:ma-models}を縦に見比べると、本章で学んだ変更がそのまま現れています。
Q/KV が等しい（MHA）のも、FFN が GELU なのも、GPT-2 だけです。
2024 年以降のモデルはすべて GQA かそれに代わる仕組みを持ち、FFN はゲート付きです。
表に列はありませんが、正規化は、GPT-2 の LayerNorm を除いて、どのモデルも RMSNorm で、位置の情報は GPT-2 以外すべて RoPE です。
2025〜2026 年に公開された 100B（1,000 億）を超える主なモデルは、すべて MoE です。
密なモデルが標準なのは、おおむね 30B（300 億）程度までです。

いくつか、本章で扱わなかった工夫も表に現れています。
DeepSeek-V3 の \textbf{MLA}（multi-head latent attention）は、キーとバリューを 512 次元の小さなベクトルに圧縮してキャッシュし、使うときに展開する方法です。
GQA とは別のやり方で、KV キャッシュを小さくします\cite{liu2024deepseekv3}。
gpt-oss や Gemma は、一部の層で「直近の百数十〜数千トークンだけを見る」\textbf{スライディングウィンドウ注意}を使い、長い文脈での計算とキャッシュを抑えています。
Qwen3 や Gemma 4 は、内積の前にクエリとキーを RMSNorm で正規化する \textbf{QK-norm} を入れて、注意のスコアが大きくなりすぎるのを防いでいます。

なお、\tabref{tab:ma-models}は 1 行に収めるために、細部を省いています。
Gemma 4 12B の「16/8」は直近だけを見る層の値で、6 層に 1 層ある、文脈全体を見る層では、キー・バリューのヘッドは 1 個です。
DeepSeek-V3 は、最初の 3 層だけ、FFN が MoE ではなく密です。
Qwen3-8B の「約 128K」は技術報告の値で、設定ファイルに書かれた位置の数は 40,960 です。
設定ファイルを開いて表と違う数字を見つけたら、こうした層ごとの違いや数え方の違いを疑ってください。

\subsection*{2026 年のモデルは、何層・何ヘッドか}

では、2026 年に公開されたモデルは、何個のヘッドを持つ注意を、何層重ねているのでしょうか。
\tabref{tab:ma-heads}に、2024〜2026 年の公開モデルの層の数 $L$、クエリのヘッドの数 $H$、キー・バリューのヘッドの数 $H_{\mathrm{kv}}$、1 ヘッドの次元 $\dhead$、隠れ次元 $d$、エキスパートの数を、本書のミニGPT と並べました。
\figref{fig:ma-heads}は、そのうちの層の数とヘッドの数を棒の長さで比べたものです。

\begin{table}[tbp]
\centering
\caption{本書のミニGPT・GPT-2 と、2024〜2026 年の公開モデルの層とヘッドの数（設定ファイルとモデルの公開ページ、技術報告の値\cite{dubey2024llama3,yang2025qwen3,liu2024deepseekv3}。「—」は手元の資料で確かめていない値）。「†」は MLA のモデルで、キャッシュするのがヘッドに分ける前の圧縮したベクトルなので、GQA のようなキー・バリューのヘッドの数がないもの。「‡」は 1 つの数で書けないもの（DeepSeek-V3 の $\dhead$ は、クエリとキーが 128＋64、バリューが 128）。$H$・$H_{\mathrm{kv}}$・$\dhead$ は、混成のモデルでは完全な注意の層の値。エキスパートは「総数／1 トークンが選ぶ数」で、「＋$s$」は全トークンが通る共有エキスパートの数。ミニGPT は第 10 章の基盤モデル（\file{base.pt}）、現代版ミニGPT は\secref{sec:ma-assemble}のもの。}
\label{tab:ma-heads}
\scriptsize
\setlength{\tabcolsep}{2pt}
\begin{tabular}{@{}lrrrrrl>{\raggedright\arraybackslash}p{44mm}@{}}
\toprule
モデル（公開） & $L$ & $H$ & $H_{\mathrm{kv}}$ & $\dhead$ & $d$ & エキスパート & 注意の構成 \\
\midrule
ミニGPT（本書・第 10 章） & 4 & 4 & 4 & 48 & 192 & なし（密） & すべて完全な注意（MHA） \\
現代版ミニGPT（本章） & 4 & 4 & 2 & 32 & 128 & なし（密） & すべて完全な注意（GQA） \\
GPT-2 124M（2019） & 12 & 12 & 12 & 64 & 768 & なし（密） & すべて完全な注意（MHA） \\
\midrule
Llama 3.1 8B（2024） & 32 & 32 & 8 & 128 & 4096 & なし（密） & すべて完全な注意（GQA） \\
Llama 3.1 405B（2024） & 126 & 128 & 8 & 128 & 16384 & なし（密） & すべて完全な注意（GQA） \\
DeepSeek-V3（2024） & 61 & 128 & † & ‡ & 7168 & 256／8＋1 & すべて MLA \\
Qwen3-8B（2025） & 36 & 32 & 8 & 128 & 4096 & なし（密） & すべて完全な注意（GQA、QK-norm） \\
Qwen3-235B-A22B（2025） & 94 & 64 & 4 & — & 4096 & 128／8 & すべて完全な注意（GQA） \\
gpt-oss-120b（2025） & 36 & 64 & 8 & 64 & 2880 & 128／4 & 1 層おきに、直近 128 トークンだけを見る注意と完全な注意 \\
\midrule
Qwen3.5-9B（2026） & 32 & 16 & 4 & 256 & 4096 & なし（密） & 3 層の線形注意（Gated DeltaNet）に 1 層の完全な注意 \\
Gemma 4 12B（2026） & 48 & 16 & 1 & 512 & 3840 & なし（密） & 5 層の窓 1,024 の注意に 1 層の完全な注意（窓の層は $H_{\mathrm{kv}}=8$、$\dhead=256$） \\
DeepSeek-V4-Pro（2026） & 61 & 128 & — & 512 & 7168 & 384／6＋1 & 圧縮した疎な注意の混成（CSA と HCA） \\
Kimi K3（2026） & 93 & 96 & † & — & 7168 & 896／16＋2 & 69 層の線形注意（Kimi Delta Attention）と 24 層のゲート付き MLA \\
Qwen3.8-27B（2026） & 64 & 24 & 4 & 256 & 5120 & なし（密） & 3 層の線形注意（Gated DeltaNet）に 1 層のゲート付き注意、の 16 回くり返し \\
Qwen3.8-2.4T-A95B（2026） & 92 & 64 & 4 & — & 8192 & 512／10＋1 & — \\
GLM-5（2026） & 78 & 64 & † & — & 6144 & 256／8＋1 & MLA と疎な注意（DSA）。層の数は設定ファイルの値（論文の本文は 80） \\
\bottomrule
\end{tabular}
\end{table}

\begin{figure}[tbp]
\centering
\begin{tikzpicture}
\begin{axis}[xbar stacked,width=0.50\linewidth,height=78mm,bar width=3.2mm,
  y dir=reverse,ytick={0,...,12},yticklabel style={font=\scriptsize},
  yticklabels={ミニGPT（本書）,現代版ミニGPT,GPT-2 124M,Llama 3.1 8B,Qwen3-8B,gpt-oss-120b,Qwen3.5-9B,Gemma 4 12B,DeepSeek-V3,Qwen3.8-27B,Kimi K3,Qwen3-235B-A22B,Llama 3.1 405B},
  xmin=0,xmax=135,enlarge y limits=0.05,xlabel={層の数 $L$},grid=none,xmajorgrids,
  legend style={font=\tiny,at={(0.98,0.97)},anchor=north east},legend cell align=left,
  title={\footnotesize (a) 層の数と、注意の種類},title style={yshift=-1mm}]
\addplot[fill=bookblue!75,draw=none] coordinates {(4,0) (4,1) (12,2) (32,3) (36,4) (18,5) (8,6) (8,7) (61,8) (16,9) (24,10) (94,11) (126,12)};
\addlegendentry{完全な注意}
\addplot[fill=bookorange!80,draw=none] coordinates {(0,0) (0,1) (0,2) (0,3) (0,4) (18,5) (0,6) (40,7) (0,8) (0,9) (0,10) (0,11) (0,12)};
\addlegendentry{直近だけを見る注意}
\addplot[fill=bookgreen!70,draw=none] coordinates {(0,0) (0,1) (0,2) (0,3) (0,4) (0,5) (24,6) (0,7) (0,8) (48,9) (69,10) (0,11) (0,12)};
\addlegendentry{線形注意}
\end{axis}
\end{tikzpicture}\hfill
\begin{tikzpicture}
\begin{axis}[xbar,width=0.40\linewidth,height=78mm,bar width=3.2mm,
  y dir=reverse,ytick={0,...,12},yticklabels={},
  xmin=0,xmax=150,enlarge y limits=0.05,xlabel={クエリのヘッドの数 $H$},grid=none,xmajorgrids,
  nodes near coords,nodes near coords align={horizontal},nodes near coords style={font=\tiny},
  title={\footnotesize (b) 1 層のヘッドの数},title style={yshift=-1mm}]
\addplot[fill=bookaccent!70,draw=none] coordinates {(4,0) (4,1) (12,2) (32,3) (32,4) (64,5) (16,6) (16,7) (128,8) (24,9) (96,10) (64,11) (128,12)};
\end{axis}
\end{tikzpicture}
\caption{本書のミニGPT と公開モデルの、層の数とヘッドの数（\tabref{tab:ma-heads}の値）。(a) 棒の長さが層の数で、色は層の種類（青: 文脈全体を見る完全な注意、橙: 直近だけを見る注意、緑: 線形注意）。色ごとの層の数を合計して並べたもので、層の並び順ではない（たとえば Qwen3.5-9B は、線形注意 3 層と完全な注意 1 層のくり返し）。(b) 1 層のクエリのヘッドの数（混成のモデルは完全な注意の層の値）。上の 2 本が本書のモデル。表のうち DeepSeek-V4-Pro と GLM-5（疎な注意の混成で、3 色に分けられない）、Qwen3.8-2.4T-A95B（注意の構成を確かめていない）の 3 個は省いた。}
\label{fig:ma-heads}
\end{figure}

\figref{fig:ma-heads}(a) の上の 2 本が、本書のモデルです。
本書のミニGPT は 4 層・4 ヘッドで、注意のヘッドは全部で $4\times4=16$ 個です。
Llama 3.1 405B は 126 層・128 ヘッドなので、$126\times128=16{,}128$ 個で、約 1,000 倍あります。
表の 2026 年のモデルは、32〜93 層に、1 層あたり 16〜128 個のクエリのヘッドを持っています。
キー・バリューのヘッドは 4〜8 個がふつうで（Gemma 4 の完全な注意の層は 1 個。MLA のモデルには、この数がありません）、クエリのヘッドよりずっと少なくなっています（GQA）。

もう一つ目立つのは、(a) の橙と緑の部分です。
2025〜2026 年のモデルには、すべての層で文脈全体を見るわけではないものが増えました。
直近の 128〜1,024 トークンだけを見る層や、線形注意の層を混ぜ、文脈全体を見る完全な注意の層を、2〜6 層に 1 層にしています（gpt-oss は 2 層に 1 層、Qwen3.5 と Qwen3.8 は 4 層に 1 層、Gemma 4 は 6 層に 1 層）。
線形注意は、過去のトークンの情報を、決まった大きさの状態にまとめながら読み進める注意で、計算量が文脈の長さ $T$ に比例します（\secref{sec:ma-beyond}）。
その層は、文脈が長くなってもキャッシュが増えません。
\figref{fig:ma-kv-mem}で Qwen3.5-9B のキャッシュが小さかったのは、このためです。

\tabref{tab:ma-attn}に、本章と本節に出てきた、注意とキャッシュの工夫を並べました。
どれを使うかは、各社で選択が分かれています。
Llama や Qwen の密なモデルは GQA を、DeepSeek・Kimi・GLM の大きな MoE は MLA を使い、2026 年には線形注意や疎な注意との混成も増えました。
どれが最善かは、まだ決着していません。

\begin{table}[tbp]
\centering
\caption{注意と KV キャッシュの工夫の比較（事実は本章の本文と、各モデルの公開された設定・技術報告による）。「キャッシュ」の列は、1 トークン・1 層あたりに増えるキャッシュが何に比例するか。}
\label{tab:ma-attn}
\footnotesize
\setlength{\tabcolsep}{3pt}
\begin{tabular}{@{}>{\raggedright\arraybackslash}p{18mm}>{\raggedright\arraybackslash}p{27mm}>{\raggedright\arraybackslash}p{17mm}>{\raggedright\arraybackslash}p{19mm}>{\raggedright\arraybackslash}p{18mm}>{\raggedright\arraybackslash}p{19mm}@{}}
\toprule
方法 & 仕組みの要点 & キャッシュ & 長所 & 短所・費用 & 使っている例 \\
\midrule
MHA & ヘッドごとにキーとバリューを持つ & $2H\dhead$ & どのヘッドも専用のキーとバリューを使える & キャッシュが最も大きい & GPT-2、Llama 2 7B、本書のミニGPT \\
GQA & $H_{\mathrm{kv}}$ 組のキーとバリューを、クエリのヘッドで分け合う & $2H_{\mathrm{kv}}\dhead$ & MHA に近い品質で、MQA に近い速さ\cite{ainslie2023gqa} & $H_{\mathrm{kv}}$ を選ぶ必要がある & Llama 3、Qwen3、gpt-oss、Gemma 4 \\
MQA & キーとバリューを 1 組だけ持つ & $2\dhead$ & キャッシュが最小 & 性能が下がることがある & 演習 4 \\
MLA & キーとバリューを低い次元のベクトルに圧縮して保存し、使うときに展開する & 圧縮した次元 & GQA とは別のやり方で小さくできる & 展開の計算が加わり、仕組みが複雑 & DeepSeek-V3、Kimi K2、GLM-5 \\
直近だけを見る注意 & 直近の $w$ トークンだけを見る（スライディングウィンドウ） & 窓の層は $w$ トークン分で頭打ち & 長い文脈で計算とキャッシュを抑える & 窓の外は直接見えないので、完全な注意の層と組み合わせる & gpt-oss（$w=128$）、Gemma 4（$w=1{,}024$） \\
線形注意との混成 & 一部の層を、過去を決まった大きさの状態にまとめる線形注意にする & 線形注意の層は増えない & 計算量が $T$ に比例 & 完全な注意の層も残す（3:1 など） & Qwen3.5、Qwen3.8、Kimi K3 \\
\bottomrule
\end{tabular}
\end{table}


\subsection*{設定からパラメータ数とメモリを計算する}

第 10 章でミニGPT のパラメータを数えたのと同じ方法で、実在のモデルも数えられます。
現代の標準形の 1 層には、注意の 4 つの射影と、SwiGLU の 3 枚の行列と、RMSNorm が 2 つあります。

\snippet{設定からパラメータ数と KV キャッシュを計算する}{code:ma-count}{ch11/param_count.py}{count}

3 行目が注意の射影で、$\mW_q$ と $\mW_o$ は $d\times H\dhead$、$\mW_k$ と $\mW_v$ は $d\times H_{\mathrm{kv}}\dhead$ です。
関数は $\dhead$ を引数 \code{dh} として別に受け取ります。
本章のコードでは $\dhead=d/H$ でしたが、実際のモデルでは $\dhead$ を別に決めることがあるからです。
たとえば gpt-oss-120b は $d=2880$、$H=64$ で $d/H=45$ ですが、$\dhead=64$ です（\tabref{tab:ma-heads}）。
KV キャッシュの式\eqref{eq:ma-kv-mem}にも、この $\dhead$ を入れます。
GQA の効果は、この $H_{\mathrm{kv}}$ の項に現れます。
4 行目が SwiGLU の $3d\,\dff$ です。
7 行目は埋め込みで、重みを共有しなければ入力側と出力側の 2 つ分になります。
MoE なら、FFN がエキスパートの数（と共有エキスパートの数）だけ並び、ルータ（$dE$、5 行目）が加わります（8 行目）。
アクティブパラメータ（9 行目）では、エキスパートの数の代わりに選ぶ数 $k$ を掛けます。
最後の \code{+ d} は最終の RMSNorm です。
\code{kv_bytes_per_token} は式\eqref{eq:ma-kv-mem}そのものです。

\outputfile{ch11/out/param_count.txt}

関数に入れた設定の値（$V$、$d$、$L$、$H$、$H_{\mathrm{kv}}$、$\dhead$、中間次元など）は、\file{ch11/param_count.py} の後半の辞書 \code{models} に並べてあり、表にあるモデルは\tabref{tab:ma-models}・\tabref{tab:ma-heads}と同じ値です。
表にない Qwen3-30B-A3B は、$V=151{,}936$、$d=2048$、$L=48$、$H=32$、$H_{\mathrm{kv}}=4$、$\dhead=128$、1 個のエキスパートの中間次元 768、128 個のエキスパートから 8 個を選ぶ設定です。

計算結果は、公開されているパラメータ数とよく一致しています。
Llama 3 8B は 8.03B、Mixtral 8x7B は総数 46.70B・アクティブ 12.88B で、名前の「8x7B」から想像する 56B より少ないことが分かります。
エキスパートになっているのは FFN だけで、注意や埋め込みは 8 個に増えていないからです。
Qwen3-30B-A3B は総数 30.53B・アクティブ 3.35B で、モデルカードの「総 30.5B・アクティブ 3.3B」と一致します（このモデルは埋め込みと出力の重みを共有しない設定です）\cite{yang2025qwen3}。
Qwen3-8B の計算値 8.19B も公開値と一致しますが、QK-norm の小さなパラメータは数えていません。

右の 2 列が KV キャッシュです。
Llama 2 7B と Llama 3 8B を比べると、パラメータ数はほぼ同じなのに、キャッシュは 512 KiB と 128 KiB で 4 倍違います。
これが GQA（$H_{\mathrm{kv}}$ が 32 → 8）の効果です。
Llama 3 70B は $L=80$ と層が多いので、1 トークンあたり 320 KiB になります。

\begin{practice}{推論に必要なメモリの見積もり}
推論に必要な GPU のメモリは、おおよそ「重み ＋ KV キャッシュ」です。
重みは総パラメータ数 × 1 つの数のバイト数（BF16 なら 2 バイト）、KV キャッシュは式\eqref{eq:ma-kv-mem} × 文脈長 × 同時に処理する系列の数です。
たとえば Llama 3 8B を BF16 で動かすと重みが約 16 GB、文脈 8,192 トークンの会話を 8 本同時に処理すると KV キャッシュが 8 GiB です。
MoE では、計算はアクティブパラメータ分でも、\textbf{メモリには総パラメータ分}が必要なことに注意してください。
量子化で重みを小さくする方法は\chapref{ch:distillation}、推論サーバの工夫は\chapref{ch:eval-serving}で扱います。
\end{practice}

この節で分かったのは、公開モデルの設定ファイルの数個の値から、パラメータ数・アクティブパラメータ・KV キャッシュの大きさを自分で計算できることです。

\section{長い文脈、マルチモーダル、思考するモデル}\label{sec:ma-beyond}
\laterread{読み物です。飛ばしても、次の章は読めます。思考するモデルの学習を扱う\chapref{ch:rl-llm}の前に、戻ってくれば足ります。}

ここからは手を動かす部分を終えて、読み物になります。
アーキテクチャの周辺で起きている大きな動きを 3 つ、概観だけしておきます。名前が出てくる工夫は、今は「そういうものがある」と知っておけば十分です。
この節の問いは、「本章の部品の上に、現代の LLM はさらに何を積み重ねているのか」です。

\textbf{長い文脈}。
文脈長は、GPT-2 の 1,024 から、2024〜2025 年の 128K（約 13 万）、2026 年には 256K〜100 万トークンへと伸びました。
それを可能にした要因の一つが RoPE です。
位置の表がないので、学習の後半に長い文書で追加の学習をすれば、文脈を伸ばせます。
そのとき、RoPE の $\theta$ を大きくしたり、YaRN~\cite{peng2024yarn} のように周波数ごとに角度の伸ばし方を変えたりして、ゆっくり回る組に学習で見たことのない角度が現れないように調整します（\secref{sec:ma-rope}）。
Llama 3 も、短い文脈で大部分を学習した後、最後の段階で 128K まで伸ばしています\cite{dubey2024llama3}。
もう一つの要因は、全トークンを見る注意の $T^2$ の費用を抑えることです。
スライディングウィンドウ注意と通常の注意を交互に重ねる方法や、計算量が $T$ に比例する「線形注意」の層を混ぜる方法が使われています（\tabref{tab:ma-heads}、\tabref{tab:ma-attn}）。

\textbf{マルチモーダル}。
画像や音声を扱うモデルも、基本の考え方は「何でもトークンの列にして、同じ Transformer に入れる」です。
画像の場合は、画像を 14×14 画素や 16×16 画素の小さな正方形（パッチ）に切り分け、各パッチを線形変換や小さな画像用のネットワークで $d$ 次元のベクトルにします。
こうしてできた「画像のトークン」を、文章のトークンの埋め込みと同じ列に並べれば、注意は画像と文章の間でも情報をやり取りできます。
Gemma 3、Qwen3.5、Llama 4 などは、こうして画像を直接扱います。

\textbf{思考（推論）するモデル}。
2024 年以降、答えの前に長い「考える過程」を生成してから答えるモデルが広まりました。
本書の語彙にも、考える過程の始まりを示す特殊トークン \code{<|think|>} を用意してあります（\chapref{ch:tokenizer}）。
アーキテクチャ自体は本章のものと同じで、違いは\textbf{学習のさせ方}にあります。
正解を確かめられる問題で強化学習を行うと、モデルは長く考えるほど正答率が上がるように振る舞いを変えていきます\cite{guo2025deepseekr1}。
これは\chapref{ch:rl-llm}の主題です。
生成するトークンが長くなるので、本章の KV キャッシュや GQA の重要性はさらに増しています。

\begin{table}[tbp]
\centering
\caption{本章の現代版ミニGPT と、実際の LLM の規模の対比。右の列の層・ヘッドの数は\tabref{tab:ma-heads}、KV キャッシュは\figref{fig:ma-kv-mem}の範囲。}
\label{tab:ma-scale}
\small
\begin{tabular}{@{}lp{34mm}p{58mm}@{}}
\toprule
項目 & 本章（現代版ミニGPT） & 実際の LLM（例） \\
\midrule
パラメータ & 約 125 万 & 80 億（Llama 3 8B）〜 6,710 億（DeepSeek-V3）。2026 年には 2 兆を超えるもの（Kimi K3 など）もある \\
層・次元 & 4 層、$d=128$ & 32〜126 層、$d=2880$〜16384（\tabref{tab:ma-heads}） \\
ヘッド Q/KV & 4/2 & 16/4〜128/8 など \\
MoE & なし（\secref{sec:ma-moe}の別の小さなモデルで 4 エキスパート、top-1） & 128〜896 エキスパートから 4〜16 個など \\
文脈長 & 128 & 128K〜100 万 \\
KV キャッシュ & 1 トークンあたり 2 KiB（32 ビットで保存） & 1 トークンあたり 32〜512 KiB 程度（\figref{fig:ma-kv-mem}） \\
学習 & CPU で約 8 分、約 160 万トークン & GPU 数千〜1.6 万枚で数か月、約 15 兆トークン \\
\bottomrule
\end{tabular}
\end{table}

\tabref{tab:ma-scale}は、本章で作ったモデルと実際の LLM の規模を並べたものです。
パラメータ数は 4〜6 桁違いますが、部品の種類と組み方は同じです。
本章の \file{modern_gpt.py} の設定の数字を表の右の列に置き換えれば、そのまま Llama 3 と同じ形のモデルになります。
足りないのは、それを学習させるための計算機とデータだけです。

この節で分かったのは、長い文脈・マルチモーダル・思考の能力が、本章の部品の上に「学習のさせ方」と「入力の作り方」を重ねて実現されていることです。

\begin{summary}
\item 現代の LLM は GPT-2 と同じ骨組み（埋め込み → ブロック × $L$ → 出力）で、ブロックの部品が入れ替わったもの。
\item RMSNorm は平均を引かずに 2 乗平均平方根で割る正規化。軽く、ずらし $\bm\beta$ を持たない。
\item RoPE はクエリとキーを位置に比例した角度だけ回し、内積に位置が「位置の差」の形でだけ入るようにする（内積はクエリとキーの中身にもよる）。位置の表がないので文脈を伸ばしやすい。
\item SwiGLU は、ゲートの枝が作る倍率を中身の枝に成分ごとに掛ける FFN。中間次元を約 $\tfrac83 d$ にすればパラメータ数は GPT-2 型とほぼ同じ。
\item KV キャッシュは過去のキーとバリューを保存して、生成の順伝播を新しいトークンの分だけにする。生成が長いとき、同じ前置きを何度も使うときによく効く。代わりにメモリを 1 トークンあたり $2LH_{\mathrm{kv}}\dhead b$ バイト使う。
\item GQA はキーとバリューのヘッドを共有して、KV キャッシュを $H_{\mathrm{kv}}/H$ 倍に縮める。
\item MoE はルータで各トークンを一部のエキスパートにだけ送り、1 トークンの計算量を増やさずにパラメータを増やす。偏りを防ぐ負荷分散が必要。
\item 公開モデルの設定の数個の値から、パラメータ数・アクティブパラメータ・KV キャッシュを計算できる。2026 年のモデルは 32〜93 層・1 層 16〜128 ヘッドで、線形注意などとの混成が増えている。
\end{summary}

\section*{次の章へ}

本章で、現代の LLM の標準的な部品は、ひととおりそろいました。
現代版ミニGPT の設定の数字を大きくすれば、Llama 3 や Qwen3 と同じ形のモデルになります。
では、なぜ本書ではそれを学習させないのでしょうか。
本章の小さなモデルですら、CPU では数分かかりました。
80 億パラメータのモデルを 15 兆トークンで学習させるのに、どれだけの計算が必要なのでしょうか。
次の\chapref{ch:gpu}では、LLM の計算量を数え、なぜ GPU が必要なのか、GPU から最大の性能を引き出すにはどうすればよいのかを考えます。

\section*{さらに学ぶために}

\begin{itemize}
\item Touvron ら「LLaMA」\cite{touvron2023llama}：RMSNorm・RoPE・SwiGLU をまとめた「現代の標準形」の出発点。短く読みやすい。
\item Su ら「RoFormer」\cite{su2024roformer}：RoPE の原論文。回転による相対位置の導出が丁寧に書かれている。
\item Shazeer「GLU Variants Improve Transformer」\cite{shazeer2020glu}：SwiGLU を含むゲート付き FFN の比較。5 ページの短い論文。
\item Ainslie ら「GQA」\cite{ainslie2023gqa}：MHA と MQA の中間としての GQA と、既存のモデルを GQA に変換する方法。
\item Fedus ら「Switch Transformers」\cite{fedus2022switch}：top-1 の MoE と負荷分散損失。MoE の実践的な入門に最適。
\item DeepSeek-AI「DeepSeek-V3 Technical Report」\cite{liu2024deepseekv3}：MLA、共有エキスパート、補助損失なしの負荷分散など、現代の大規模 MoE の設計がすべて載っている。
\item Grattafiori ら「The Llama 3 Herd of Models」\cite{dubey2024llama3}：密なモデルの設計・学習・文脈の延長の詳細な記録。
\item Raschka『Build a Large Language Model (From Scratch)』\cite{raschka2024build}：GPT-2 型から Llama 型への書き換えを、コードで追える。
\end{itemize}

\begin{exercises}
\exfirst[どちらも、必読の節の式に数を入れて確かめる問題です。]{2 と 5}
\item （手計算）$\vx=(2,-2,1,-1)$、$\bm\gamma=(1,1,1,1)$、$\epsilon=0$ のとき、RMSNorm と LayerNorm（$\bm\beta=\vzero$）の出力をそれぞれ求めてください。平均が 0 のベクトルでは、2 つはどんな関係になりますか。
\item （手計算）2 次元の RoPE で $\omega=0.5$、$\vq=(1,0)$、$\vk=(0,1)$ とします。$m=4,\,n=2$ と $m=10,\,n=8$ のスコアを計算し、一致することを確かめてください。
\item （計算）$d=4096$ で、GPT-2 型の MLP（$\dff=4d$）と同じパラメータ数になる SwiGLU の $\dff$ を求めてください。Llama 2 7B の 11008 と比べてどうですか。
\item （コード）\file{ch11/kv_cache.py} で \code{n_kv_head} を 4 から 1（MQA）に変えて、速度を測り直してください。また、\code{KVCache} に保存されたテンソルの形を表示して、メモリが何分の 1 になったか確かめてください。
\item （計算）Llama 3.1 405B（$L=126$、$H_{\mathrm{kv}}=8$、$\dhead=128$）で、BF16 の KV キャッシュは 1 トークンあたり何 KiB ですか。文脈 131,072 トークンの会話 1 本では何 GiB になりますか。 同じ計算を gpt-oss-120b（$L=36$、$H_{\mathrm{kv}}=8$、$\dhead=64$）でも行い、比べてください。
\item （コード）\file{ch11/moe_experiment.py} で top-1 を top-2 に変えて学習し、損失とエキスパートの割合がどう変わるか調べてください。1 トークンあたりの FFN の計算量は何倍になりますか。
\item （考察）MoE では、アクティブパラメータが小さくても、総パラメータ分のメモリが必要です。手元の GPU 1 枚で動かしたいとき、密な 80 億パラメータのモデルと、総 300 億・アクティブ 30 億の MoE の、どちらが向いているでしょうか。速さとメモリの両面から考えてください。
\end{exercises}
